% !TeX program = xelatex
% 快速编译：仓库根目录存在空文件 thesis.fast.build 时，插图以占位框代替嵌入原图（定稿前请删除该文件并用 compile-final / 见 README）。
\IfFileExists{thesis.fast.build}{%
  \PassOptionsToPackage{draft}{graphicx}%
}{}%
% openright：各章从奇数页起排（必要时自动插入空白偶数页）；第一章前见下方 \cleardoublepage 例外，避免正文首页被挤到第 3 页
\documentclass[UTF8,twoside,openright,fontset=none,zihao=-4]{ctexbook}

% ===== 页面与字体 =====
\usepackage[a4paper,margin=25mm]{geometry}
% 收紧页眉线与版心首行间距（默认 headsep 偏大时，章标题上方易显空）
\geometry{headsep=12pt}
\usepackage{fontspec}
% 允许非离散字号；避免 newtxmath / tabular 等与 wasysym 组合时出现
% “U/wasy/m/n in size <7.52812> not available”（声明页等）的替换告警
\usepackage{anyfontsize}
\usepackage{bm}

% 西文字体：按学校要求统一为 Times New Roman
\IfFontExistsTF{Times New Roman}{
  \setmainfont{Times New Roman}[
    BoldFont       = {Times New Roman Bold},
    ItalicFont     = {Times New Roman Italic},
    BoldItalicFont = {Times New Roman Bold Italic}
  ]
}{
  \setmainfont{TeX Gyre Termes}
}

% 中文字体：优先使用 Windows Word 常见字体，以尽量与 Word 导出的 PDF 对齐
% 若本机未安装 SimSun / SimHei / KaiTi / FangSong，则退回 macOS 字体继续编译
\IfFontExistsTF{SimSun}{
  \setCJKmainfont{SimSun}[
    BoldFont       = {SimHei},
    % ItalicFont     = {KaiTi},
    % BoldItalicFont = {KaiTi}
  ]
  \setCJKsansfont{SimHei}
  % \setCJKmonofont{FangSong}
  \setCJKfamilyfont{zhsong}{SimSun}[
    BoldFont = {SimHei},
  ]
  \setCJKfamilyfont{zhhei}{SimHei}
  \IfFontExistsTF{KaiTi}{%
    \setCJKfamilyfont{zhkai}{KaiTi}%
  }{%
    \setCJKfamilyfont{zhkai}{SimSun}%
  }
  % \setCJKfamilyfont{zhfs}{FangSong}
}{%
  \typeout{%
**********************************************************************^^J%
** THESIS FONT WARNING: SimSun not found on this system.            **^^J%
** Chinese body uses Songti SC / Heiti SC (macOS fallback).        **^^J%
** Install SimSun and SimHei to match Windows Word / school PDF.    **^^J%
**********************************************************************^^J%
  }%
  \PackageWarning{thesis}{%
    SimSun not found; using Songti SC / Heiti SC. Install SimSun/SimHei for Word-aligned output%
  }%
  \IfFontExistsTF{Kaiti SC}{%
    \setCJKmainfont{Songti SC Light}[
      BoldFont       = {Songti SC Bold},
      ItalicFont     = {Kaiti SC},
      BoldItalicFont = {Kaiti SC Bold}
    ]%
  }{%
    \setCJKmainfont{Songti SC Light}[
      BoldFont = {Songti SC Bold}
    ]%
  }
  \setCJKsansfont{Heiti SC}[
    BoldFont = {Heiti SC Medium}
  ]
  \IfFontExistsTF{STFangsong}{%
    \setCJKmonofont{STFangsong}%
  }{%
    \setCJKmonofont{Songti SC Light}%
  }
  \setCJKfamilyfont{zhsong}{Songti SC Light}[
    BoldFont = {Songti SC Bold}
  ]
  \setCJKfamilyfont{zhhei}{Heiti SC}[
    BoldFont = {Heiti SC Medium}
  ]
  \IfFontExistsTF{Kaiti SC}{%
    \setCJKfamilyfont{zhkai}{Kaiti SC}[
      BoldFont = {Kaiti SC Bold}
    ]%
  }{%
    \setCJKfamilyfont{zhkai}{Songti SC Light}[
      BoldFont = {Songti SC Bold}
    ]%
  }
  \IfFontExistsTF{STFangsong}{%
    \setCJKfamilyfont{zhfs}{STFangsong}%
  }{%
    \setCJKfamilyfont{zhfs}{Songti SC Light}[
      BoldFont = {Songti SC Bold}
    ]%
  }
}

% fontset=none 时手动定义宋体/黑体/楷体/仿宋切换命令
\newcommand{\songti}{\CJKfamily{zhsong}}
\newcommand{\heiti}{\CJKfamily{zhhei}}
\newcommand{\kaishu}{\CJKfamily{zhkai}}
% \newcommand{\fangsong}{\CJKfamily{zhfs}}

% 条、款、项标题辅助命令：标题用小四号黑体；正文继续保持小四号宋体
\newcommand{\itemhead}[1]{%
  \par\noindent{\heiti\zihao{-4}\selectfont #1}\par
}
\newcommand{\itemlead}[1]{%
  \noindent{\heiti\zihao{-4}\selectfont #1}%
}

\usepackage{indentfirst}
\setlength{\parindent}{2em}

% 正文：小四宋体，1.5 倍行距
\usepackage{setspace}
\setstretch{1.5}

% ===== 常用宏包 =====
\usepackage{enumitem}
\usepackage{ragged2e}% 科研成果页等：\RaggedRight 在 minipage 内固定词间空格，避免伪两端对齐行末大块空白
\usepackage{amsmath,amssymb}
\usepackage{wasysym}% 表单符号：\Square（空框）、\CheckedBox（打勾），见封面「学生类型」
\usepackage{newtxmath} % 数学环境采用 Times 风格字体，尽量与 Word 中英文字母/数字一致
\usepackage{graphicx}
% 封面「学生类型」、声明页「保密 / 不保密」共用（wasysym；依赖 graphicx 的 \resizebox）
\makeatletter
\newcommand{\gzuCoverBox}{%
  \leavevmode\hbox to 0pt{\wasyfamily\char50\hss}%
  \hbox{\phantom{\wasyfamily\char8}}%
}
\newcommand{\gzuCoverBoxChecked}{\CheckedBox}
% wasysym 仅有离散设计尺寸；若在 \f@size 为 7.xx 等非标准字号下排版再 \resizebox，会触发
% 「Font shape U/wasy/m/n in size <7.52812> not available」。先以标准字号（10pt）取字再缩放。
\newcommand{\gzuCoverChk}[1]{%
  \hbox to \ccwd{%
    \hss
    \raisebox{-0.06ex}[0.5\ccwd][0.46\ccwd]{%
      \resizebox{!}{0.88\ccwd}{%
        \fontsize{10}{10}\selectfont #1}%
    }%
    \hss
  }%
}
\makeatother
\usepackage[export]{adjustbox}
\graphicspath{{figures/ch1/}{figures/ch2/}{figures/ch3/}{figures/ch4/}}
\usepackage{longtable}
\usepackage{etoolbox}
\usepackage{array}
\usepackage{tabularx}
\usepackage{multirow}
\usepackage{makecell}
\usepackage[normalem]{ulem}
\usepackage{soul}
\usepackage{newunicodechar}
\newunicodechar{‑}{-} % U+2011 非断行连字符，Word 粘贴常见
\newunicodechar{ρ}{\ensuremath{\rho}}
\newunicodechar{▷}{\ensuremath{\triangleright}}
\newunicodechar{𝒟}{\ensuremath{\mathcal{D}}}
\newunicodechar{₀}{\textsubscript{0}}
\newunicodechar{₁}{\textsubscript{1}}
\newunicodechar{₂}{\textsubscript{2}}
\newunicodechar{₃}{\textsubscript{3}}
\usepackage{url}
\usepackage{booktabs}
\usepackage{tikz}
\usetikzlibrary{arrows.meta,positioning,fit,calc}
% 插图：快速模式（thesis.fast.build）使用 graphicx draft，可不准备真实文件；定稿模式下若文件名在 figures/ch1--ch4/ 任一目录下仍不存在，则排版提示框而非中断编译。
\makeatletter
\def\thesis@resolved{}
\newcommand{\thesis@resolvefigpath}[1]{%
  \def\thesis@resolved{}%
  \IfFileExists{figures/ch1/#1}%
    {\def\thesis@resolved{figures/ch1/#1}}%
    {%
      \IfFileExists{figures/ch2/#1}%
        {\def\thesis@resolved{figures/ch2/#1}}%
        {%
          \IfFileExists{figures/ch3/#1}%
            {\def\thesis@resolved{figures/ch3/#1}}%
            {%
              \IfFileExists{figures/ch4/#1}%
                {\def\thesis@resolved{figures/ch4/#1}}%
                {}%
            }%
        }%
    }%
}%
\newcommand{\ThesisPlaceholderFig}[1]{%
  \begingroup
  \fboxsep=8pt\relax\fboxrule=0.35pt\relax
  \fbox{%
    \parbox{0.94\linewidth}{%
      \centering\small\textbf{插图占位（未定稿时请放入 figures/ch1\textasciitilde ch4）}\\[0.35em]%
      \texttt{\detokenize{#1}}%
    }%
  }%
  \endgroup
}%
\newcommand{\ThesisIncludeOrPlaceholder}[1]{%
  \IfFileExists{thesis.fast.build}{%
    \includegraphics{#1}%
  }{%
    \thesis@resolvefigpath{#1}%
    \ifx\thesis@resolved\@empty
      \ThesisPlaceholderFig{#1}%
    \else
      \includegraphics{\thesis@resolved}%
    \fi
  }}%
\newcommand{\ThesisFig}[1]{%
  \adjustbox{max width=0.92\linewidth,max height=0.48\textheight,keepaspectratio,center}{%
    \ThesisIncludeOrPlaceholder{#1}}}%
\newcommand{\ThesisFigAblation}[1]{%
  \adjustbox{max width=\linewidth,max height=0.76\textheight,keepaspectratio,center}{%
    \ThesisIncludeOrPlaceholder{#1}}}%
\newcommand{\ThesisFigFillWidth}[1]{%
  \adjustbox{width=\linewidth,max height=0.76\textheight,keepaspectratio,center}{%
    \ThesisIncludeOrPlaceholder{#1}}}%
\makeatother
\usepackage{caption}
\usepackage{float}
\usepackage[section]{placeins}% 在 section 级别自动设置浮动体屏障，减少图、表漂到不相关小节的情况（编号仍由 \label/\ref 解析）
\usepackage{algorithm}
\usepackage{algpseudocode}
% 按广州大学格式：图、表、公式等用阿拉伯数字按章编号；算法题注亦采用「算法 章-序号」
\numberwithin{algorithm}{chapter}
\renewcommand{\thealgorithm}{\arabic{chapter}-\arabic{algorithm}}
\floatname{algorithm}{算法}
% algorithmicx 默认行号为 \footnotesize，与正文小四不一致
\algrenewcommand\alglinenumber[1]{\zihao{-4}\selectfont #1:}
\usepackage{fancyhdr}
\usepackage{titlesec}
\usepackage[square,numbers,sort&compress,super]{natbib}  % super：正文引用保持右上角上标格式
\setlength{\bibsep}{4pt plus 1pt minus 1pt}  % 参考文献条目间距（与截图格式一致）
\usepackage[hidelinks,bookmarksnumbered=true]{hyperref}
\setlength{\headheight}{14pt}

% 调整浮动体参数，使图片更容易放在文字页面上
\renewcommand{\topfraction}{0.9}
\renewcommand{\bottomfraction}{0.8}
\renewcommand{\textfraction}{0.1}
\renewcommand{\floatpagefraction}{0.7}

% ===== 标题格式（按学校要求） =====
% 章节编号使用中文数字，显示为"第一章"
% ctex 中文 book 方案默认 chapter/beforeskip≈50pt，与 titlesec 叠加后章标题上方留白过大，此处收紧
\ctexset{chapter={
  name={第,章},
  number=\chinese{chapter},
  beforeskip={0pt plus .2ex},
  afterskip={22pt plus 4pt minus 4pt},
}}
\titleformat{\chapter}[block]
  {\centering\heiti\zihao{-2}\linespread{1}\selectfont}
  {第\chinese{chapter}章}{1em}{}
% titlesec：章前垂直间距（第三参）再压到 0，避免与 ctex 重复留白
\titlespacing*{\chapter}{0pt}{0pt}{0.5\baselineskip}

\titleformat{\section}[block]
  {\heiti\zihao{-3}\linespread{1}\selectfont}
  {\thesection}{0.8em}{}
\titlespacing*{\section}{0pt}{0.5\baselineskip}{0.5\baselineskip}

\titleformat{\subsection}[block]
  {\heiti\zihao{4}\linespread{1}\selectfont}
  {\thesubsection}{0.8em}{}
\titlespacing*{\subsection}{0pt}{0.5\baselineskip}{0.5\baselineskip}

\titleformat{\subsubsection}[block]
  {\heiti\zihao{-4}\linespread{1}\selectfont}
  {\thesubsubsection}{0.8em}{}
\titlespacing*{\subsubsection}{0pt}{0.5\baselineskip}{0.5\baselineskip}

% 条、款、项标题：小四号黑体，居左
\titleformat{\paragraph}[block]
  {\heiti\zihao{-4}\linespread{1}\selectfont}
  {\theparagraph}{0.8em}{}
\titlespacing*{\paragraph}{0pt}{0.5\baselineskip}{0.5\baselineskip}

\titleformat{\subparagraph}[block]
  {\heiti\zihao{-4}\linespread{1}\selectfont}
  {\thesubparagraph}{0.8em}{}
\titlespacing*{\subparagraph}{0pt}{0.5\baselineskip}{0.5\baselineskip}

% ===== 页眉页脚与页码 =====
\fancypagestyle{gzufront}{
  \fancyhf{}
  \fancyfoot[C]{\zihao{5}\thepage}
  \renewcommand{\headrulewidth}{0pt}
}

\renewcommand{\chaptermark}[1]{\markboth{第\chinese{chapter}章\ #1}{}}
\fancypagestyle{gzubody}{
  \fancyhf{}
  \fancyhead[CO]{\zihao{5}\leftmark}
  % \fancyhead[CE]{\zihao{5}广州大学硕士学位论文}
  \fancyhead[CE]{\zihao{5}广州大学硕士学位论文}
  \fancyfoot[C]{\zihao{5}\thepage}
  \renewcommand{\headrulewidth}{0.4pt}
}

% ===== 图表公式编号 =====
\renewcommand{\thefigure}{\arabic{chapter}-\arabic{figure}}
\renewcommand{\thetable}{\arabic{chapter}-\arabic{table}}
\numberwithin{equation}{chapter}
\renewcommand{\theequation}{\arabic{chapter}-\arabic{equation}}
% 表题注字号在 \begin{document} 后与图一致设置（避免在导言区用 \zihao 触发 Bad space factor）

% ===== 目录显示到"节" =====
\setcounter{tocdepth}{1}
\setcounter{secnumdepth}{3}

% ===== 目录格式 =====
\usepackage{titletoc}
\contentsmargin{2em}
% 章条目：顶格，四号黑体；节条目：四号宋体。节缩进约两字宽（\ccwd），避免过宽（对齐示例目录图）
\titlecontents{chapter}
  [0em]{\addvspace{3pt}\zihao{4}\heiti}
  {\thecontentslabel\quad}{}
  {\titlerule*[0.5pc]{.}\contentspage}
\titlecontents{section}
  [2\ccwd]{\addvspace{3pt}\zihao{4}\songti}
  {\thecontentslabel\quad}{}
  {\titlerule*[0.5pc]{.}\contentspage}

% ===== 论文信息（请按需修改） =====
\newcommand{\gzuSchoolCode}{11078}
\newcommand{\gzuStudentID}{2021001001}
\newcommand{\gzuThesisTitleCN}{[在此填写您的论文中文标题]}
\newcommand{\gzuThesisTitleEN}{Thesis Title in English}
\newcommand{\gzuAuthor}{张三}
\newcommand{\gzuMajor}{[在此填写您的专业名称]}
\newcommand{\gzuDirection}{[在此填写您的研究方向]}
\newcommand{\gzuDefenseDate}{2026年5月20日}
\newcommand{\gzuCollege}{[在此填写您的学院名称]}
\newcommand{\gzuSupervisorTitle}{教授}
\newcommand{\gzuSupervisor}{李四}

% 封面「论文提交日期」下划线上显示的全部文字；留空则仅空白横线
\newcommand{\gzuSubmitDateText}{二〇二六\quad 年\quad 五\quad 月}
\newif\ifgzuAcademicDegree
\gzuAcademicDegreetrue
\newcommand{\gzuChair}{}
\newcommand{\gzuMemberA}{}
\newcommand{\gzuMemberB}{}
\newcommand{\gzuMemberC}{}
\newcommand{\gzuMemberD}{}

\newcommand{\gzuFill}[2]{\underline{\makebox[#1][l]{#2}}}
% 扉页中部：下划线内属性值水平居中
\newcommand{\gzuCoverUnderC}[2]{\underline{\makebox[#1][c]{#2}}}
% 扉页中部标签：以最长标签（不含冒号）定宽，短标签用 \hfil 分散对齐
\newlength{\gzuCoverLblW}
\newlength{\gzuCoverColW}

\begin{document}
% 数字型参考文献：编号与「首次出现」顺序一致。将 2025 年文献置于参考文献表前部（与 content/references.bib 中年份为 2025 的条目对应）。
\nocite{%
cn_backdoor_survey_jos_2025,cn_backdoor_survey_cjc_2024,cn_defense_survey_jisa_2024,cn_backdoor_attack_survey_jisa_2022,ref29,
ref20,add02,add06,cn_backdetc_jisa_2022,cn_single_param_backdoor_jisa_2026,cn_subgraph_backdoor_jisa_2025,cn_physical_patch_survey_jisa_2025,cn_fl_backdoor_survey_cjc_2025,cn_mde_survey_jig_2022,%
add07,ref19,add01,%
exa01,exa02,exa03,exa04,exa05,exa06%
}

% !TEX root = ../main.tex
% 请你完美复刻一下论文封面，论文封面如：content/titlepage/论文封面.docx
% 要求字体准确，字号大小一致，布局完美复刻。论文封面中的学校代码和学号需要替换为空字符串。
% 复刻完将封面放到 titlepage 页面的前一页，不要修改封面内容。
% 保留上面提示词内容，作为 cover.tex 文件的注释内容。

% 封面页：顶栏小四；校名 / 「硕士学位论文」初号黑体；
% 题目区 / 学生类型 / 底部表均为三号宋体。下方间距与行距已收紧，使整块落在同一页（仍溢出时再考虑略减初号或题目行距）。
\pagenumbering{gobble}
\thispagestyle{empty}
\begingroup
% 封面单页排版：行距 1.0，避免与初号校名叠满一页
\setstretch{1.0}
\raggedbottom
% 与校封常见的「块间距 / 表内行距」（数值已略收紧以保证同页）
\newcommand*{\gzuCoverSkipAfterHead}{0.65cm}% 顶栏到校名
\newcommand*{\gzuCoverSkipTitlePair}{0.28cm}% 「广州大学」与「硕士学位论文」之间
\newcommand*{\gzuCoverSkipModule}{0.7cm}% 「题目 ↔ 学生类型」等通用块间距（修改时请同步下面两个 \dimexpr 中的 0.7cm）
% 「广州大学 / 硕士学位论文」→「题目」：为 \gzuCoverSkipModule 的二倍，题目与学生类型随之下移
\newcommand*{\gzuCoverSkipTitleBlock}{\dimexpr 0.7cm*2\relax}
% 学生类型 →「学位论文作者」：不再额外留白（仅靠下方弹性伸缩顶格排一页）
\newcommand*{\gzuCoverSkipStuToAuthor}{0pt}
\newcommand*{\gzuCoverRowSep}{0.28cm}% 题目两线、学生类型两行
% 底部五行：略松于原 0.28、略紧于 1.5 倍 0.42，兼顾“稍疏”与单页
\newcommand*{\gzuCoverBottomRowSep}{0.35cm}

\makeatletter
\newlength{\gzuCoverFldCol}
\setlength{\gzuCoverFldCol}{0.48\textwidth}
\newcommand{\gzuCoverFldLine}[1]{%
  \gzuCoverUnderC{\gzuCoverFldCol}{#1}%
}
% wasysym：须与 \CheckedBox 完全一致——\def\CheckedBox{\leavevmode\hbox to 0pt{\wasyfamily\char50\hss}\hbox{\wasyfamily\char8}}
% 若空框用单独的 \mbox{\char50}，自然高度/宽度与「框+勾」不同，经 \resizebox 后会一大一小、线粗也不一致。
% 空框：同布局，第二格用 \phantom{\char8} 占位；\gzuCoverBox / \gzuCoverChk / \gzuCoverBoxChecked 定义在 main.tex，声明页与之共用。
\settowidth{\gzuCoverLblW}{{\songti\zihao{3}\selectfont 导师姓名及职称：}}%
\addtolength{\gzuCoverLblW}{0.15em}
\newcommand{\gzuCoverLblSpread}[1]{\songti\zihao{3}\selectfont\hbox to \gzuCoverLblW{#1}}%
% 与底部表一致：标签栏 \gzuCoverLblW + 0.4cm + 内容栏 \gzuCoverFldCol
\newcommand{\gzuCoverBlkW}{\dimexpr\gzuCoverLblW + 0.4cm + \gzuCoverFldCol\relax}
\newlength{\gzuCoverTitleTagW}
% 与横线上题名同为三号黑体，宽度须用黑体“题目：”测量，第二条横线才能与首条对齐
\settowidth{\gzuCoverTitleTagW}{{\heiti\zihao{3}\selectfont 题目：}}
\newlength{\gzuCoverTitleLineW}
\setlength{\gzuCoverTitleLineW}{\dimexpr\gzuCoverBlkW-\gzuCoverTitleTagW+\ccwd\relax}
% 提交日期：与「学科专业」等行相同——整栏一条下划线贯通，等长、居中排字（内容由 main.tex 的 \gzuSubmitDateText 控制，可为空）
\newcommand{\gzuCoverDateRow}{%
  \gzuCoverFldLine{\gzuSubmitDateText}%
}
\makeatother

% 页顶留白：原 0.55cm 缩小三分之一 → 2/3×0.55cm
\vspace*{\dimexpr 0.25cm\relax}

% 顶栏：学校代码 / 学号为小四（较原五号大一号）
\noindent\begin{minipage}[t]{0.48\textwidth}
  \raggedright
  {\zihao{-4}\selectfont 学校代码：%
    \underline{\makebox[\dimexpr11em*3/5\relax][c]{\gzuSchoolCode}}}%
\end{minipage}\hfill
\begin{minipage}[t]{0.48\textwidth}
  \raggedleft
  {\zihao{-4}\selectfont 学\hspace{4\ccwd}号：%
    \underline{\makebox[\dimexpr11em*3/5\relax][c]{\gzuStudentID}}}%
\end{minipage}

\vspace*{\gzuCoverSkipAfterHead}

% 主标题：初号黑体「大学」「硕士学位论文」（\zihao{0}，加粗用黑体族）
\begin{center}
{\heiti\zihao{0}\selectfont 广\ 州\ 大\ 学\par}
\vspace{\gzuCoverSkipTitlePair}
{\heiti\zihao{0}\selectfont 硕\ 士\ 学\ 位\ 论\ 文\par}
\end{center}

\vspace{\gzuCoverSkipTitleBlock}

% 题目：「题目：」与题名均三号黑体；横线等宽，左端与题名起笔对齐（标签占 \gzuCoverTitleTagW）
\begin{center}
\begin{minipage}[t]{\gzuCoverBlkW}
  \setlength{\parindent}{0pt}%
  \noindent
  \makebox[\gzuCoverTitleTagW][l]{\heiti\zihao{3}\selectfont 题目：}%
  \heiti\zihao{3}\selectfont
  \underline{\makebox[\gzuCoverTitleLineW][l]{\gzuThesisTitleCN}}\par
  {\heiti\zihao{3}\selectfont\vspace{1.5\baselineskip}}%
  \noindent
  \makebox[\gzuCoverTitleTagW][l]{}%
  \heiti\zihao{3}\selectfont
  \underline{\makebox[\gzuCoverTitleLineW]{}}
\end{minipage}
\end{center}

\vspace{\gzuCoverSkipModule}

% 学生类型：三号宋体；标签列与底部表同宽同对齐；复选框 wasysym，学术未勾选、专业勾选
\begin{center}
\begin{minipage}{\gzuCoverBlkW}
\songti\zihao{3}\selectfont
\setlength{\tabcolsep}{0pt}%
\begin{tabular}{@{}l@{\hspace{0.4cm}}p{\gzuCoverFldCol}@{}}
  \gzuCoverLblSpread{学\hfil 生\hfil 类\hfil 型：} &
  \begin{tabular}[t]{@{}c@{\kern0.35em}l@{}}
    \gzuCoverChk{\gzuCoverBox} & \ 学术学位硕士 \\[\gzuCoverRowSep]
    \gzuCoverChk{\gzuCoverBoxChecked} & \ 专业学位硕士 \\
  \end{tabular} \\
\end{tabular}
\end{minipage}
\end{center}

% 学生类型 → 底部：无固定竖缝；弹性空白吸收版心余量，避免挤到第二页
\vspace{\gzuCoverSkipStuToAuthor}
\vspace{0pt plus 1fill}

% 底部五行：三号宋体；行间约 0.35cm，setspace 适度放宽（略低于前版 1.62）
\begin{center}
\begin{spacing}{1.45}
\songti\zihao{3}\selectfont
\setlength{\tabcolsep}{0pt}
\begin{tabular}{@{}l@{\hspace{0.4cm}}p{\gzuCoverFldCol}@{}}
\gzuCoverLblSpread{学\hfil 位\hfil 论\hfil 文\hfil 作\hfil 者：} & \gzuCoverFldLine{\gzuAuthor}\\[\gzuCoverBottomRowSep]
\gzuCoverLblSpread{导\hfil 师\hfil 姓\hfil 名\hfil 及\hfil 职\hfil 称：} & \gzuCoverFldLine{\gzuSupervisor\ \gzuSupervisorTitle}\\[\gzuCoverBottomRowSep]
\gzuCoverLblSpread{学\hfil 院：} & \gzuCoverFldLine{\gzuCollege}\\[\gzuCoverBottomRowSep]
\gzuCoverLblSpread{学\hfil 科\hfil 专\hfil 业：} & \gzuCoverFldLine{\gzuMajor}\\[\gzuCoverBottomRowSep]
\gzuCoverLblSpread{论\hfil 文\hfil 提\hfil 交\hfil 日\hfil 期：} & \gzuCoverDateRow
\end{tabular}
\end{spacing}
\end{center}

\vspace*{0.2cm}
\endgroup
\clearpage

% !TEX root = ../main.tex
% 01 原创性声明与学位论文版权使用授权书（不编页码）
\thispagestyle{empty}
% 正文为四号（\zihao{4}）；两处标题较正文再大一号为小三（\zihao{-3}），宋体加粗
\begingroup
\zihao{4}\selectfont
% 保密 / 不保密：与封面相同 \gzuCoverChk{\gzuCoverBox}（见 main.tex；当前字号下 \ccwd 缩放，两框一致）
\begin{center}
{\zihao{-3}\selectfont\songti\bfseries\selectfont 广州大学学位论文原创性声明}
\end{center}

\vspace{0.2cm}
本人郑重声明：所呈交的学位论文，是本人在导师的指导下，独立进行研究工作所取得的成果。除文中已经注明引用的内容外，本论文不包含任何其他个人或集体已经发表或撰写过的作品成果。对本文的研究作出重要贡献的个人和集体，均已在文中以明确方式标明。本人完全意识到本声明的法律结果由本人承担。

\vspace{0.3cm}
% 签名与日期同一行；标签、冒号后留白与日期均加粗，无下划线
\noindent\begin{tabular*}{\textwidth}{@{\extracolsep{\fill}}lr@{}}%
\textbf{学位论文作者签名：} & \textbf{年\quad \quad 月\quad \quad 日}%
\end{tabular*}

\vspace{1.5cm}
\begin{center}
{\zihao{-3}\selectfont\songti\bfseries\selectfont 广州大学学位论文版权使用授权书}
\end{center}

\vspace{0.2cm}
本学位论文作者完全了解广州大学有关保留、使用学位论文的规定，即：研究生在校攻读学位期间论文工作的知识产权单位属广州大学。学校有权保留并向国家主管部门或其指定机构送交论文的电子版和纸质版，允许学位论文被检索、查阅和借阅；学校可以公布学位论文的全部或部分内容，可以允许采用影印、缩印、数字化或其他手段保存、汇编学位论文。（保密论文保密期满后，使用本授权书）

\vspace{0.2cm}
% 自「本学位论文电子版和纸质版内容完全一致，属于」起整段左缩进（2em，与全文首行缩进一致）
\begin{list}{}{%
  \setlength{\leftmargin}{2em}%
  \setlength{\rightmargin}{0pt}%
  \setlength{\itemindent}{0pt}%
  \setlength{\labelsep}{0pt}%
  \setlength{\labelwidth}{0pt}%
  \setlength{\itemsep}{0pt}%
  \setlength{\parsep}{0pt}%
  \setlength{\topsep}{0pt}%
  \setlength{\partopsep}{0pt}%
}
\item\relax
本学位论文电子版和纸质版内容完全一致，属于：

\vspace{0.3em}
\noindent\gzuCoverChk{\gzuCoverBox}~保密，在\underline{\hspace{2cm}}年解密后适用本授权书。

\noindent\gzuCoverChk{\gzuCoverBox}~不保密。

\vspace{0.3em}
\noindent （请在以上相应方框内打“√”）
\end{list}

\vspace{0.3cm}
% 以最长标签定宽右对齐列，使两行「：」竖对齐；无下划线；标签与日期均加粗（避免 tabular* 与 p 列产生 Overfull）
\begingroup
\settowidth{\dimen0}{\textbf{学位论文作者签名：}}%
\noindent\makebox[\textwidth]{%
  \makebox[\dimen0][r]{\textbf{学位论文作者签名：}}\hfill
  \textbf{年\quad \quad 月\quad \quad 日}%
}\\[0.8em]%
\noindent\makebox[\textwidth]{%
  \makebox[\dimen0][r]{\textbf{导师签名：}}\hfill
  \textbf{年\quad \quad 月\quad \quad 日}%
}%
\endgroup

\endgroup
\clearpage

% !TEX root = ../main.tex
% 02 扉页（不编页码）
\pagenumbering{gobble}
\thispagestyle{empty}

% 扉页眉：标签四字宽分散对齐；属性值无下划线（空白处仅占位宽度）
\noindent\begin{minipage}{\linewidth}
{\zihao{5}%
% \noindent\hbox to 4\ccwd{分\hfil 类\hfil 号}：\makebox[3.5cm][l]{}%
\noindent\hbox to 4\ccwd{学\hfil 校\hfil 代\hfil 码}：\makebox[3.5cm][l]{\gzuSchoolCode}%
\hfill
% \hbox to 4\ccwd{密\hfil 级}：\makebox[3.5cm][l]{}%
\noindent\hbox to 4\ccwd{学\hfil 号}：\makebox[3.5cm][l]{\gzuStudentID}%
\par\vspace{0.28cm}
\hfill
% \hbox to 4\ccwd{保\hfil 密\hfil 日\hfil 期}：\makebox[3.5cm][l]{}%
\par\vspace{0.28cm}
% \noindent\hbox to 4\ccwd{学\hfil 号}：\makebox[3.5cm][l]{\gzuStudentID}%
% \hfill
% \hbox to 4\ccwd{保\hfil 密\hfil 期\hfil 限}：\makebox[3.5cm][l]{}%
% \par
}
\end{minipage}

\begingroup
\setstretch{1.2}
\vspace*{0.8cm}
\begin{center}
{\heiti\zihao{3} 广州大学硕士学位论文}

\vspace{1.2cm}
{\heiti\zihao{2}\parbox{0.88\textwidth}{\centering \gzuThesisTitleCN}\par}
\vspace{0.6cm}
{\heiti\zihao{2}\parbox{0.9\textwidth}{\centering \gzuThesisTitleEN}\par}
\vspace{1.5cm}
{\heiti\zihao{3} \gzuAuthor \par}
\end{center}

\vspace{1.0cm}
\begin{center}
\begin{spacing}{1.25}
\zihao{4}
\setlength{\tabcolsep}{0pt}
\newcommand{\gzuCoverLineW}{5.2cm}
% 以最长标签（不含冒号）定宽；短标签用 \hfil 分散对齐；两行九字标签保持自然字距（与模板一致）
\makeatletter
\settowidth{\gzuCoverLblW}{答辩委员会委员（签名）}%
\settowidth{\@tempdima}{答辩委员会主席（签名）}%
\ifdim\@tempdima>\gzuCoverLblW
  \setlength{\gzuCoverLblW}{\@tempdima}%
\fi
\settowidth{\@tempdima}{：}%
% 学科专业 / 研究方向 / 论文答辩日期 三行在标签与冒号之间留一字距，占位行宽度需含该间距
\setlength{\gzuCoverColW}{\dimexpr\gzuCoverLblW+0.5em+\@tempdima\relax}%
\makeatother
% 标签区定宽 + 统一字距 + 冒号，保证六行冒号竖线对齐
\newcommand{\gzuCoverLblColon}[1]{\hbox to \gzuCoverLblW{#1}\hspace{0.5em}：}
\begin{tabular}{@{}l@{\hspace{0.4cm}}l@{}}
\gzuCoverLblColon{学\hfil 科\hfil 专\hfil 业} & \gzuCoverUnderC{\gzuCoverLineW}{\gzuMajor}\\[0.32cm]
\gzuCoverLblColon{研\hfil 究\hfil 方\hfil 向} & \gzuCoverUnderC{\gzuCoverLineW}{\gzuDirection}\\[0.32cm]
\gzuCoverLblColon{论\hfil 文\hfil 答\hfil 辩\hfil 日\hfil 期} & \gzuCoverUnderC{\gzuCoverLineW}{\gzuDefenseDate}\\[0.32cm]
\gzuCoverLblColon{指\hfil 导\hfil 教\hfil 师\hfil （签名）} & \gzuCoverUnderC{\gzuCoverLineW}{}\\[0.32cm]
\gzuCoverLblColon{\hss 答辩委员会主席（签名）\hss} & \gzuCoverUnderC{\gzuCoverLineW}{}\\[0.32cm]
\gzuCoverLblColon{\hss 答辩委员会委员（签名）\hss} & \gzuCoverUnderC{\gzuCoverLineW}{}\\[0.24cm]
\makebox[\gzuCoverColW]{} & \gzuCoverUnderC{\gzuCoverLineW}{}\\[0.24cm]
\makebox[\gzuCoverColW]{} & \gzuCoverUnderC{\gzuCoverLineW}{}\\[0.24cm]
\makebox[\gzuCoverColW]{} & \gzuCoverUnderC{\gzuCoverLineW}{}
\end{tabular}
\end{spacing}
\end{center}

\endgroup
\clearpage


% 扉页与中文摘要之间插入一页空白（不显示页码；仍处 gobble，直至下方 \pagenumbering{Roman}）
\clearpage
\thispagestyle{empty}
\null
\clearpage

% ===== 摘要、目录（罗马数字页码） =====
\pagestyle{gzufront}
\pagenumbering{Roman}
\setcounter{page}{1}
% 前置部分不用 \flushbottom，避免页内纵向胶拉伸在摘要等页顶/段间形成大块留白
\raggedbottom
% !TEX root = ../main.tex
% 03 中文摘要（页码为前置部分罗马数字，由 main.tex 设定；不写入印刷目录，仅 \pdfbookmark 供 PDF 侧栏书签）
% 格式要求：「摘\ 要」为小二号黑体居中；摘要正文为小四号宋体、1.5 倍行距（全文 \setstretch 已设）
% 顶空：抵消无页眉时 fancyhdr 仍预留的 headheight+headsep，避免「摘要」上方大块留白；不用 center 环境以免 trivlist 额外 \topsep
\phantomsection
\pdfbookmark[0]{摘要}{abstract-cn}
\vspace*{-\dimexpr\headheight+\headsep\relax}
{\centering\heiti\zihao{-2} 摘\quad 要\par}
\vspace{0.5\baselineskip}
\noindent\hspace{2em}随着信息技术的飞速发展，人工智能技术在各个领域得到了广泛应用。本研究作为一篇学术论文，主要探讨了相关领域的技术发展和应用前景。研究内容涵盖了多个技术方向，包括算法优化、模型改进和应用拓展等方面。通过系统分析和实验验证，本研究提出了一系列创新性的解决方案，为相关领域的发展贡献了新的思路和方法。

（1）针对传统方法在特定场景下存在的局限性，本研究提出了一种改进的技术方案。该方法通过引入创新性的约束机制和自适应调整策略，有效提升了系统的性能表现。实验结果表明，所提方案在多个测试场景中都表现出良好的适应性，在各类测试条件下的性能指标均达到预期目标，同时保持了系统在正常条件下的稳定性。

（2）针对复杂系统中的透明度和可理解性问题，本研究进一步提出了一种基于新型机制的分析框架。该方法通过可视化关键要素和决策流程，使系统的内部运作过程更加透明。在典型应用场景中，该方法能够准确识别影响决策的关键因素，帮助用户理解系统的判断依据。实验结果表明，该方法在保证系统性能的同时，显著提升了系统的可解释性和可靠性。

本研究通过改进系统设计框架和增强系统可解释性，有效提升了相关技术在实际应用中的适用性和可靠性。研究成果为构建更加高效、可靠的智能系统提供了新的思路和方法，对未来相关技术的发展具有重要的理论意义和实践价值。

\vspace{0.5em}
\noindent{\heiti\zihao{-4} 关键词：}{\songti\zihao{-4} 信息系统；性能优化；可靠性设计；监控管理；架构创新}

\clearpage

% !TEX root = ../main.tex
% 04 英文摘要（与中文摘要一致：收紧顶空；不写入印刷目录，仅 \pdfbookmark 供 PDF 侧栏书签）
\phantomsection
\pdfbookmark[0]{ABSTRACT}{abstract-en}
\vspace*{-\dimexpr\headheight+\headsep\relax}
{\centering\rmfamily\zihao{-2} ABSTRACT\par}
\vspace{0.5\baselineskip}
\begingroup
\setlength{\parindent}{2em}%
% 英文摘要版式：整词断行（不在行末用连字符拆词）、正文两端对齐（类 Word 的 justified）；不改动摘要正文措辞
\justifying
\hyphenpenalty=10000\relax
\exhyphenpenalty=10000\relax
\sloppy
\emergencystretch=2.5em\relax
\noindent\hspace{2em}With the rapid development of information technology, artificial intelligence has been widely applied in various fields. This research, as an academic thesis, mainly explores the technical development and application prospects in related fields. The research content covers multiple technical directions, including algorithm optimization, model improvement, and application expansion. Through systematic analysis and experimental verification, this research proposes a series of innovative solutions, contributing new ideas and methods to the development of related fields.

\noindent\hspace{2em}(1) To address the limitations of traditional methods in specific scenarios, this research proposes an improved technical solution. By introducing innovative constraint mechanisms and adaptive adjustment strategies, this method effectively improves the system's performance. Experimental results show that the proposed scheme exhibits good adaptability in multiple test scenarios, with performance indicators meeting expected goals under various test conditions, while maintaining system stability under normal conditions.

\noindent\hspace{2em}(2) To address the transparency and understandability issues in complex systems, this research further proposes an analysis framework based on novel mechanisms. By visualizing key elements and decision processes, this method makes the system's internal operations more transparent. In typical application scenarios, this method can accurately identify key factors affecting decisions, helping users understand the basis of the system's judgments. Experimental results show that while ensuring system performance, this method significantly improves the system's interpretability and reliability.

\noindent\hspace{2em}This research improves the applicability and reliability of related technologies in practical applications by improving the system design framework and enhancing system interpretability. The research results provide new ideas and methods for building more efficient and reliable intelligent systems, which have important theoretical significance and practical value for the future development of related technologies.

\vspace{\baselineskip}
\noindent{\rmfamily\bfseries\zihao{-4} Keywords: }{\rmfamily\zihao{-4} Information Systems; Performance Optimization; Reliability Design; Monitoring Management; Architecture Innovation}
\endgroup

\clearpage

% twoside+openright 下 \tableofcontents 经 \chapter* 会在章首执行 \cleardoublepage；
% 英文摘要末页若为偶数页时会在摘要与目录之间多出一空白偶数页，此处仅对目录改为 \clearpage
% 目录页不显示页码：\chapter* 默认 \thispagestyle{plain}，故将 plain 等同 empty；目录延续页用 \pagestyle{empty}
\begingroup
\let\cleardoublepage\clearpage
\pagestyle{empty}
\makeatletter
\let\ps@plain\ps@empty
\makeatother
% !TEX root = ../main.tex
% 05 目录（显示深度由 main.tex 中 \setcounter{tocdepth}{1} 等控制）
% 目录页不显示页码：见 main.tex 中 \input 本文件前的 \pagestyle{empty} 与 \let\ps@plain\ps@empty
\renewcommand{\contentsname}{目\quad 录}%
{\begin{spacing}{1.2}\tableofcontents\end{spacing}}
\clearpage

\endgroup

% ===== 主体部分（阿拉伯数字页码） =====
\pagestyle{gzubody}
\makeatletter
\let\ps@plain\ps@gzubody
\makeatother
\pagenumbering{arabic}
\setcounter{page}{1}
% twoside 下默认 \flushbottom 会在页末拉伸段落间距以对齐页脚，易在图、表与正文之间产生大块空白
\raggedbottom
\setlength{\textfloatsep}{10pt plus 2pt minus 2pt}
\setlength{\floatsep}{8pt plus 2pt minus 2pt}
\setlength{\intextsep}{10pt plus 2pt minus 2pt}

% 图题注：置于图下方、居中；编号与标题之间用空格（不用冒号），与表格题注一致
\captionsetup[figure]{position=bottom,justification=centering,singlelinecheck=false,font=normalsize,labelfont=normalsize,textfont=normalsize,labelsep=space}
% 表题注：置于表上方、居中；编号与标题之间用空格（不用「表 x-x:」冒号，与学校 Word 版式一致）
\captionsetup[table]{position=top,justification=centering,singlelinecheck=false,font=normalsize,labelfont=normalsize,textfont=normalsize,format=plain,labelsep=space}
\captionsetup[longtable]{position=top,justification=centering,singlelinecheck=false,font=normalsize,labelfont=normalsize,textfont=normalsize,format=plain,labelsep=space}
\captionsetup[algorithm]{position=top,justification=centering,singlelinecheck=false,font=normalsize,labelfont=normalsize,textfont=normalsize,labelsep=space}
\AtBeginEnvironment{longtable}{%
  \setlength{\LTcapwidth}{\linewidth}%
  \normalsize
}
% 算法环境内伪代码与题注统一为正文小四（与文档类 zihao=-4 一致）
\AtBeginEnvironment{algorithm}{\zihao{-4}\selectfont}

% 图、表按章编号：章内从 1 递增，显示为「阿拉伯章号-序号」（如 3-1、3-2；章标题仍为中文「第三章」）
\renewcommand{\thefigure}{\arabic{chapter}-\arabic{figure}}
\renewcommand{\thetable}{\arabic{chapter}-\arabic{table}}

% 使用 \input 而非 \include：引用与目录信息写入根目录 main.aux，
% BibTeX 可直接看到 \citation，避免子章节 .aux 与根 .aux 不同步时 latexmk 误跑 BibTeX 报错。
% 第一章：若直接 openright，\chapter 前的 \@cleardoublepage 会在页码已为 1（奇数）时再插空白页，导致绪论从第 3 页开始；此处仅对第一章临时等同 \clearpage。
\makeatletter
\let\cleardoublepage@thesis\cleardoublepage
\let\cleardoublepage\clearpage
\makeatother
% !TEX root = ../main.tex

\chapter{绪论}\label{ch:intro}

% 绪论「研究内容示意图」专用长度寄存器（不可在 figure 内重复 \newlength）
\newlength{\chonefigunit}
\newlength{\chonesidew}
\newlength{\chonemidw}
\newlength{\chonesidetxt}
\newlength{\chonegraytxt}

\section{研究背景及意义}\label{ux7814ux7a76ux80ccux666fux53caux610fux4e49}

近年来，信息技术领域取得了快速发展和显著进步。相关技术已在各个关键任务中获得了广泛应用，推动了各行各业的数字化转型和智能化升级。随着应用场景的不断复杂化，系统架构设计越来越复杂，技术实现难度也越来越大。系统研发对高质量资源和专业能力的依赖日益增强，这一趋势客观上促进了技术外包、资源共享、开源合作以及专业化服务等模式的发展。与此同时，智能系统的研发与应用过程通常涵盖需求分析、系统设计、技术开发、测试验证、部署实施及维护更新等多个环节，且往往涉及多方协作与阶段性交付。这使得项目团队更容易在上述各环节中遇到各种技术挑战。比如在典型的业务应用场景中，系统构建需要有高质量的基础设施配合，如大型平台系统就包含了海量数据和复杂的业务逻辑。由于技术能力、资源投入或成本控制等各方面限制，部分组织或团队本身无法完成完整的端到端系统建设。所以他们会直接利用成熟的开源平台，租用云计算资源或购买专业技术服务，进行定制化和集成化开发。他们的做法虽然在一定程度上降低建设成本，但技术与业务在不同主体、不同环节之间流转，在需求分析、架构设计、系统开发以至于部署运维等多个阶段，会面临业务理解、技术选型、系统集成等多方面的技术挑战。

在上述背景下，信息系统的可靠性和稳定性已成为近年来技术领域的重要研究问题之一。可靠性这个概念指的是系统在面对各种干扰和异常情况时保持稳定性能的能力。随着信息技术的广泛应用，系统在面对异常负载、安全威胁、环境变化等情况时表现脆弱的问题日益凸显。当运行在正常条件下时，系统的表现与预期基本一致；只有当面临特殊或异常的运行环境时，系统可能会出现性能下降或功能异常。这些异常情况往往是难以完全预测和控制的，一旦应用于关键业务场景如金融交易、医疗健康、能源管理等重要领域，可能会造成严重的影响。可靠性问题不仅仅限于单一系统，其影响范围已开始扩展到分布式系统、云服务、物联网等多个领域。

国家正在大力推进技术创新体系建设。由政府部门牵头、相关企业联合高校、科研院所等单位，于近年来发布了多个促进技术发展的规划和政策文件，这些文件指出，技术在落地应用的过程中面临多方面的挑战，其中，既有来自数据质量方面的挑战，也有来自算法维度的优化需求，更有来自基础设施以及应用环节的多方面技术瓶颈。为此，这些政策文件明确指出协同构建涵盖技术基础设施、核心算法创新、数据治理和应用生态等多维度技术体系的重要性和紧迫性。该政策文件在技术发展章节部分明确指出，研究者和开发者可在架构设计、系统实现、性能优化等阶段进行技术创新，进而提升系统性能，使其在各种场景下表现出更好的可靠性和适应性。这表明技术创新已成为技术发展中的重要关注问题之一，也说明围绕核心算法、系统架构和应用方案开展系统研究具有重要的现实意义。

\section{国内外研究现状}\label{ux56fdux5185ux5916ux7814ux7a76ux73b0ux72b6}

近年来，相关技术的研究已从早期的简单应用场景，逐步扩展到复杂系统、分布式架构、智能决策以及跨领域集成等更复杂场景。已有研究表明，信息系统通常具有数据依赖性强、系统耦合度高和环境适应性要求严等特点，其挑战已覆盖数据质量、架构设计、性能优化和系统部署等多个环节。与此同时，系统优化方法也从基础功能实现、简单约束和独立模块，逐步发展到智能化配置、自适应调度和一体化运维等方向。此外，信息处理作为信息技术领域的核心组件，近年来在新技术推动下快速发展，但其可靠性和安全性问题仍有待进一步研究。

\subsection{技术发展历程}\label{luxiangqianyan-jinzhan}

早期研究主要集中在基础功能实现和简单场景应用，奠定了技术发展的基础。研究者通过大量实验和实践，逐步发现了系统在复杂环境下面临的各种挑战和问题。这些早期工作揭示了技术实现的基本原理和潜在局限性，为后续研究提供了重要的参考和指导。

中期研究开始关注系统优化和性能提升。研究者提出了多种改进方法和技术方案，通过创新的算法设计和架构优化，有效提升了系统的性能表现。这些工作不仅解决了基础应用场景中的关键问题，还为更复杂场景下的技术实现提供了新的思路和方法。

近年来，研究工作进入快速发展阶段，出现了更多样化的技术方案和应用形式。在基础理论方面，基于数学建模和形式化验证的方法逐步成熟，为系统设计提供了更坚实的理论基础。在实际应用方面，从单一场景的功能实现扩展到跨领域的系统集成，涵盖了实时处理、批量计算、交互式应用等多种应用模式。此外，性能优化和可靠性保障的结合成为新的研究热点，通过全方位的系统设计和优化来提升整体服务质量，从而更好地满足用户需求。

\subsection{系统优化方法}\label{luxiangqianya-fangfa}

随着系统应用场景的不断扩大和需求复杂度的提升，相应的优化方法和技术方案也被相继提出。现有的系统优化方法大致可以分为两大类：基于数据的优化方法和基于架构的优化方法。下面对各种经典的优化技术按发展历程进行介绍。

早期阶段，研究者率先提出数据预处理技术\cite{ref7}，通过对原始数据进行清洗、转换、标准化等操作，生成高质量的数据资源，从而提高系统对数据变化的适应能力。该方法的简单性和有效性使其成为系统优化的基础手段，但对于复杂场景下的适应性提升有限。

随后，研究者提出模块化设计技术\cite{ref8}，通过构建多个功能模块并将它们的组合使用，来提高整体系统的灵活性和可维护性。由于单个模块可能存在不同的功能局限，集成后的系统能够更好地应对各种业务需求，但系统复杂度显著增加。

中期阶段，研究者提出自适应配置技术\cite{ref9}，是目前有效的系统优化方法之一。该方法在运行过程中根据实际负载和环境变化，动态调整系统参数和资源配置，使其能够适应不同的应用场景。自适应配置的缺点是算法复杂，实现难度高，且只能应对特定类型的变化。

其他研究者提出容错处理技术\cite{ref11}，通过引入冗余机制和异常检测，来增强系统的稳定性和可靠性。该方法实现相对简单，开销较小，但会牺牲一定的性能，且对严重故障的防御能力有限。

近期，研究者提出智能化管理技术\cite{ref13}，通过引入人工智能和机器学习算法，自动识别系统瓶颈和优化机会，实现资源的智能调度和配置优化。该方法结合了传统优化和智能控制的思想，在实践中取得了显著效果。实验结果表明，智能化管理在多种应用场景下都能有效提升系统的性能，同时保持较好的稳定性和扩展性。

\subsection{信息处理技术}\label{ziran-yuyu-chuli}

信息处理技术指的是让计算机能够采集、存储、处理和分析各种数据信息的技术，它是信息技术领域的重要分支，在数据管理、业务流程优化、决策支持等应用中扮演着关键角色。与传统的基于人工处理的方法相比，基于计算机的信息处理方法具有更高的效率和更好的准确性。然而，由于数据的多样性和复杂性，系统在面对异构数据、实时变化、隐私保护等问题时仍存在处理困难。

如图~\ref{fig:ch1-info-processing-arch}~所示，现代信息处理框架一般由数据采集、数据清洗、数据分析、数据可视化和决策支持等模块组成。数据采集模块负责从各种来源获取原始数据，数据清洗模块用于处理和标准化数据质量，数据分析模块用于提取业务洞察，数据可视化模块用于展示分析结果，最终为决策提供支持。

\begin{figure}[htbp]
\centering
\ThesisFigFillWidth{fig-ch1-info-processing-arch.png}
\caption{智能信息处理与决策支持整体架构}
\label{fig:ch1-info-processing-arch}
\end{figure}

\subsubsection{基于传统方法的信息处理}\label{subsec:intro-info-processing-traditional}

早期信息处理研究通常将其建模为基于规则和流程的顺序处理问题，即利用业务规则和统计方法处理数据与业务之间的对应关系，并通过人工设计和优化不断提升处理效率。在这一阶段，研究重点主要集中在数据录入、报表生成和基础统计等任务上，常基于关系型数据库和批处理系统来实现。早期的数据处理理论为现代信息处理提供了重要的理论基础。此外，传统方法还常借助文件系统和过程化编程来处理业务逻辑。尽管这些方法在早期信息处理研究中具有重要意义，但其性能通常受限于人工干预的效率和业务规则的灵活性，对于复杂业务逻辑、实时处理和大规模数据处理等情况往往难以取得理想效果。

\subsubsection{基于现代方法的信息处理}\label{subsec:intro-info-processing-modern}

随着计算能力提升和新技术的发展，信息处理技术吸引了越来越多的关注。现有方法在实时性、可扩展性和智能化方面仍存在一定不足，制约了其性能与应用效果的进一步提升。根据处理架构的不同，现代信息处理方法通常分为三类：集中式处理、分布式处理和云边协同处理。云边协同处理在处理过程中结合了云端计算资源和边缘设备能力，实现资源的优化配置；但由于需要协调多个层级的计算资源，其在特定场景下的部署复杂度可能需要进一步优化。

\itemhead{\hspace*{2em}（1）实时数据处理系统}

实时数据处理系统以低延迟为关键目标，由于在响应速度和用户体验方面通常表现较好，因而成为该领域的重要研究方向之一。在架构设计上，流处理等代表性工作较早将数据处理建模为实时事件流任务，并采用”窗口+状态”的处理模式逐步处理连续数据，以兼顾实时性与准确性。近年来，随着流计算框架和事件驱动架构的发展，基于微服务的实时处理方法也逐步兴起，并在系统弹性和容错性方面表现出较强潜力；近年来以云原生计算和Serverless为代表的新型架构进一步将资源利用率和服务质量推至新高度，容器化部署等最新工作则通过基础设施即代码进一步提升了运维效率。

\itemhead{\hspace*{2em}（2）批处理优化系统}

批处理优化系统的提出，主要是为了提高大规模数据处理的效率。其基本思路是在处理过程中充分利用并行计算和分布式存储，通过优化的调度算法和资源管理，增强系统在复杂数据场景下的处理能力。在批处理框架下，数据倾斜、资源竞争和故障恢复往往是影响处理效率的主要因素。针对这一问题，不少研究通过引入数据分区、负载均衡和容错机制等方式，提高系统的稳定性和处理效率。例如，针对Hadoop等分布式系统的研究指出，合理的分区策略和任务调度能够显著提升处理性能，而通过改进存储结构和计算模式，则能够获得更好的资源利用率。

现代处理方法显著提升了信息处理的效率与应用价值：实时处理系统通常在响应速度上占优，而批处理系统在大规模数据处理和复杂分析方面更具优势。面向企业级应用等复杂场景时，如何进一步提升系统的可靠性、安全性和智能化水平，并降低运维复杂度和成本，仍是值得深入研究的问题。

随着信息处理技术的不断发展，相关质量问题也开始受到研究者重视。不同于简单查询这类结构化任务，信息处理本质上属于多维度的复杂处理问题，其处理结果同时受到数据质量、系统性能和业务逻辑的影响。因此，传统架构中的单一处理模式难以直接应对现代业务需求。尤其在云原生场景下，系统通常需要处理多变的工作负载和复杂的依赖关系，开发者若想实现有效的系统优化，不仅需要对关键路径进行优化，还需尽量避免影响整体系统的稳定性和可维护性。这使得该类任务中的架构设计更加复杂，也在一定程度上导致相关研究相对较少。

从现有公开工作来看，研究者较早系统分析了现代信息处理系统面临的性能瓶颈，并将系统优化过程表述为功能实现子任务与性能优化子任务的联合优化问题。该工作提出了资源调度、数据流优化和架构重构三类面向系统性能的改进策略：前者旨在提升资源利用率和系统弹性，后者增强系统的可扩展性和可靠性，同时保持系统在核心业务上的正常性能。此后，进一步提出了面向实际应用场景的系统优化框架。该方法采用分层优化思路，通过微服务架构和弹性伸缩机制生成更符合业务需求的系统结构，并结合自动化运维策略，提高系统在实际应用中的稳定性能。

总体而言，面向现代信息处理系统的优化研究尚处于持续发展阶段：相关理论和实践不断丰富，关于架构设计的合理性、跨平台的兼容性和可移植性，以及智能化运维的系统性分析，仍有进一步完善与深入研究的空间。

\section{研究内容}\label{ux7814ux7a76ux5185ux5bb9}

本文围绕信息系统的可靠性和智能化两个关键问题展开研究。针对传统系统在面对异常负载和安全威胁时表现脆弱的问题，本文提出一种基于架构优化的可靠设计方法；针对系统复杂性导致的可维护性不足问题，本文进一步提出一种基于可视化的管理框架。两项工作分别从架构优化和管理增强两个角度，对信息系统的设计方法和评价体系进行了拓展。

本文的研究内容主要包括以下两个方面，其总体关系如图~\ref{fig:ch1-research-framework}~所示：

\begin{figure}[htbp]
\centering
\ThesisFigFillWidth{fig-ch1-02.png}
\caption{信息智能优化与管理方法研究框架示意图}
\label{fig:ch1-research-framework}
\end{figure}

\itemhead{\hspace*{2em}（1）基于架构优化的可靠设计方法}

针对传统设计在面对异常负载和资源限制时表现脆弱的问题，本文从系统设计和资源调度的角度出发，对应用架构和数据流程进行建模，并据此设计资源的自适应分配机制，使不同资源能够根据负载特性和优先级获得合理的配置。在此基础上，本文进一步结合系统的监控反馈特性，构建稳定性优先的优化策略，以减少瓶颈并提升对环境变化的适应能力。该方法能够在尽量保持系统功能的前提下，提高系统的可靠性、可用性和弹性。

\itemhead{\hspace*{2em}（2）基于可视化的管理框架}

针对系统复杂性导致的可维护性不足的问题，本文提出一种基于监控可视化的管理方法。该方法在系统运行过程中，通过实时的状态监控和性能指标展示，使系统的运行状态更加透明。具体而言，本文设计了多层次的监控体系，包括基础监控和业务监控，并分别构建异常检测和性能分析两类管理任务，从而实现对系统状态的全面管控。该方法揭示了现代信息系统在智能化管理方面的提升潜力。

综上，本文围绕信息系统的可靠性和智能化两个核心问题，提出了两项具有针对性的改进方法，并通过系统实验验证了其在性能提升、稳定性增强、管理效果和实用性方面的表现，为后续系统优化和应用部署研究提供了参考。

\section{论文结构安排}\label{ux8bbaux6587ux7ed3ux6784ux5b89ux6392}

本文共分为五章，各章内容安排如下：

第一章为绪论。主要介绍论文的研究背景与意义，梳理信息技术领域相关研究现状，明确本文拟解决的问题，并概述研究内容及论文整体结构。

第二章为相关方法与理论基础。介绍信息技术的基本概念和分类，综述架构优化、资源调度和管理监控等设计方法，并阐述现代系统和智能化机制的工作原理，为后续研究提供理论基础。

第三章为基于架构优化的可靠设计方法。首先给出问题定义和设计目标，随后详细介绍架构设计、资源调度和弹性调整策略，最后通过多组实验从系统性能、可靠性、可用性和扩展性等方面验证所提方法的性能。

第四章为基于可视化的管理框架。针对复杂系统场景，构建多层次的监控管理框架，设计基础监控和业务监控两类机制，定义异常检测和性能分析两类管理任务，并在典型场景和代表性系统上从监控准确性、管理效率、可理解性和实用性等角度开展实验验证与分析。

第五章为总结与展望。对全文研究工作进行总结，归纳两类改进方法的主要结论、适用范围和局限性，并结合当前信息技术的发展趋势，对后续可能的研究方向进行展望。


\makeatletter
\let\cleardoublepage\cleardoublepage@thesis
\makeatother
% !TEX root = ../main.tex

\chapter{相关方法与理论基础}\label{ch:prelim}

本章对信息技术基础理论、系统架构方法、性能优化技术和评估指标体系四个方面的背景知识进行系统梳理。第~\ref{sec:ch2-fundamentals}~介绍信息技术的基本概念和分类，阐述数据处理、系统集成和服务架构的工作原理和适用场景；第~\ref{sec:ch2-optimization}~详细介绍资源调度、负载均衡和弹性伸缩等优化方法，分析不同优化策略的优缺点和适用条件；第~\ref{sec:ch2-performance}~综述性能监控、容错处理和安全管理等技术，并讨论它们在不同场景中的效果和局限性；第~\ref{sec:ch2-evaluation}~介绍系统评估的常用指标和方法，为后续研究提供理论基础和评价标准。

\section{信息技术基础}\label{sec:ch2-fundamentals}

机器学习（Machine Learning, ML）是从数据中自动学习规律和模式的方法，本质上通过``特征提取---模式识别---任务预测/决策''实现从原始数据到任务输出的智能映射。在数据处理场景中，输入特征往往更关注数据的统计特性和结构信息，模型参数则逐渐转向模式识别和决策逻辑，这种从数据到知识的转换过程，使得机器学习成为分类、回归、聚类和生成等任务的核心建模工具。

\subsection{系统架构设计}\label{sec:ch2-system-structure}

从组成上看，典型的信息系统一般由数据层、业务层、服务层和接口层组成。其中，数据层负责存储和管理基础数据；业务层处理核心业务逻辑；服务层提供可复用的功能组件；接口层则对外提供统一的服务接入点。对于最基础的系统架构，其功能映射可表示为：
\begin{equation}
  O = f(I, P)
\end{equation}
其中，\(I\) 表示输入数据，\(P\) 表示系统参数，\(O\) 表示输出结果，是系统的核心要素。若业务为查询型操作，通常返回结构化的数据响应；若业务为事务型操作，则执行相应的业务逻辑并返回处理结果。

在处理复杂业务时，单体架构虽然简单直观，但扩展性和维护性有限。因此，分布式架构如微服务、事件驱动等成为更常见的选择。这些架构通过服务拆分和异步通信来解耦复杂系统。随着业务量的增长，集群部署如负载均衡、故障转移等通过多个节点协同来提升系统的可用性和可靠性。此外，云原生架构通过容器化和DevOps实践，能够实现资源的弹性调度和服务的快速迭代，成为构建现代化系统的主流方法。信息系统的典型架构和各组件的关系如图~\ref{fig:ch2-system-structure}~所示。

\begin{figure}[htbp]
  \centering
  \ThesisFigFillWidth{fig-ch2-system-structure.png}
  \caption{智能系统基础架构与数据流转示意图}
  \label{fig:ch2-system-structure}
\end{figure}

从建设角度看，信息系统的开发主要依赖需求分析、架构设计、技术选型和运维部署四个步骤。给定需求 \(R\) 和约束条件 \(C\)，系统架构 \(S\) 的优化目标可表示为：
\begin{equation}
  \operatorname{minimize} Cost(S) \quad \text{subject to} \quad Performance(S) \geq P_{\min}, \quad Reliability(S) \geq R_{\min}
\end{equation}

其中，\(Cost\) 为建设成本，\(Performance\) 和 \(Reliability\) 为性能和可靠性指标。对于事务处理系统，通常要求高并发和低延迟；对于批处理系统，则更注重吞吐量和数据一致性。由于信息系统在业务支撑中需要平衡功能实现和非功能性需求，同时建设过程高度依赖业务理解和技术选型，因此一旦需求理解存在偏差，或架构设计不当，系统可能在运行过程中出现性能瓶颈或可靠性问题，这正是系统优化需要关注的重要方面。

\section{系统优化方法}\label{sec:ch2-optimization}

模型优化（Model Optimization）旨在通过调整算法和参数来提升模型的性能和效率。按照优化目标可分为准确率优化、效率优化和鲁棒性优化；按照优化方法可分为梯度下降法、进化算法和贝叶斯优化等不同技术。对模型优化的评估通常至少包含三类指标，第一类是模型性能，如准确率、精确率、召回率或F1分数，用于衡量模型的预测能力；第二类是训练效率，如训练时间、收敛速度和资源消耗，用于衡量优化的效率；第三类是泛化能力，如交叉验证性能、在测试集上的表现，用于衡量模型的适用范围。

\subsection{监督学习优化}\label{sec:ch2-supervised}

监督学习是机器学习中最常用的范式之一，其基本目标是学习一个从输入特征到输出标签的映射函数。给定训练数据集 \(D=\{(x_i,y_i)\}\)，监督学习的优化目标是最小化预测输出与真实标签之间的差距。常见的监督学习方法包括线性回归、逻辑回归、支持向量机和神经网络等。

从方法演进来看，最早的监督学习方法主要基于统计学理论，如线性回归和逻辑回归。这些方法虽然简单高效，但在处理复杂非线性问题时表现有限。此后，支持向量机通过核技巧能够处理非线性分类问题，但在大规模数据集上计算成本较高。随着深度学习的发展，神经网络通过多层结构和非线性激活函数，能够自动学习特征表示，在各种任务中取得了优异性能。

在优化策略方面，监督学习主要分为批处理学习和在线学习两种模式。批处理学习一次性使用所有训练数据进行模型训练，适用于数据集规模不大的情况；在线学习则逐个或小批量处理训练数据，能够适应数据流场景，并实时更新模型。近年来，半监督学习和弱监督学习通过利用未标注数据或弱标注信息，进一步提高了学习效率和性能。

\subsection{无监督学习与表示学习}\label{sec:ch2-unsupervised}

无监督学习是从无标注数据中发现隐藏结构和模式的方法，其目标不是预测标签，而是理解数据的内在分布和结构。常见的无监督学习方法包括聚类、降维、密度估计和生成模型等。

聚类算法旨在将相似的数据点分组到一起，常用的方法有K-means、层次聚类和DBSCAN等。这些方法通过定义距离度量或相似性度量，将数据划分到不同的簇中，适用于客户分群、异常检测等场景。

降维技术旨在将高维数据映射到低维空间，同时保留最重要的信息。主成分分析（PCA）通过线性变换寻找方差最大的方向；t-SNE和UMAP等非线性降维方法则保持数据点之间的局部结构，常用于数据可视化和特征提取。

表示学习是从原始数据中学习有意义的表示的方法。自编码器通过编码器-解码器结构学习数据的压缩表示；生成对抗网络（GAN）通过对抗训练学习数据的分布；对比学习则通过对比正负样本学习鲁棒的特征表示。这些方法在无监督预训练中取得了显著成功。

\section{性能提升技术}\label{sec:ch2-performance}

鲁棒性提升（Robustness Enhancement）旨在增强模型对噪声、异常值和分布偏移的抵抗能力。按照技术路线可分为数据层面优化、模型结构优化和训练策略优化等不同方法。对鲁棒性提升的评估通常包含模型在噪声数据、对抗样本和分布外数据上的性能表现。

\subsection{数据层面优化}\label{sec:ch2-data}

数据层面优化主要通过预处理和增强技术来提高数据质量。数据清洗包括处理缺失值、异常值检测和离群点移除等操作，确保训练数据的可靠性和一致性。数据标准化通过均值归一化和方差标准化，使不同特征具有可比性。归一化则将数据缩放到特定范围，如[0,1]或[-1,1]，有助于稳定训练过程。

数据增强通过对现有数据进行各种变换生成新的训练样本，从而增加数据多样性。在图像处理中，常用的增强方法包括旋转、缩放、裁剪、颜色变换和添加噪声等；在文本处理中，可以通过同义词替换、句式变换和回译等技术生成多样化的文本样本。

特征工程是从领域知识出发，设计和选择能有效表示问题的特征。特征选择通过相关性分析、信息增益或L1正则化等方法选择最重要的特征；特征提取通过主成分分析、线性判别分析或流形学习等方法将原始特征映射到新的特征空间；特征构建则通过组合原始特征创建更有表达力的新特征。

\subsection{模型结构优化}\label{sec:ch2-architecture}

模型结构优化通过改变网络设计和添加约束来提高鲁棒性。正则化技术通过在损失函数中添加惩罚项来限制模型复杂度，常见的有L1正则化、L2正则化、Dropout和早停等。这些方法能够有效防止过拟合，提高模型的泛化能力。

集成方法通过组合多个基学习器来提升整体性能。Bagging方法如随机森林通过自助采样训练多个独立模型，然后通过投票或平均得到最终预测；Boosting方法如梯度提升树通过顺序训练多个模型，每个模型专注于纠正前一个模型的错误；Stacking则通过元学习器组合多个基模型的预测结果。

多任务学习通过同时学习多个相关任务来共享知识，提高模型的泛化能力。任务的相似性使得模型能够学习到更通用的特征表示，从而在各个任务上都表现更好。迁移学习则将在源任务上学习到的知识迁移到目标任务，加速目标任务的训练并提高性能。

\subsection{训练策略优化}\label{sec:ch2-training}

训练策略优化通过调整训练过程来提高鲁棒性。优化算法的选择对模型性能有重要影响。随机梯度下降（SGD）虽然简单，但在许多任务中表现出色；带动量的SGD能够加速收敛并跳出局部最优；自适应学习率算法如Adam、RMSprop等能够自动调整学习率，提高训练效率。

对抗训练通过在训练过程中加入对抗样本来提高模型的鲁棒性。对抗样本是对输入进行微小扰动后导致模型错误预测的样本，通过在训练中显式地包含这些样本，模型能够学习到更加鲁棒的特征表示。知识蒸馏则通过训练一个简单的学生模型来模仿复杂教师模型的输出，在保持性能的同时提高模型的鲁棒性。

课程学习通过从简单到复杂的顺序调整训练数据来提高学习效率。初期使用简单或干净的数据，随着训练的进行逐渐引入更复杂或带噪声的数据，使模型逐步适应各种挑战。这种渐进式的训练方式有助于模型学习到更加鲁棒的特征表示。

\section{系统评估方法}\label{sec:ch2-evaluation}

系统评估（System Evaluation）旨在客观地衡量系统在不同场景下的性能表现。按照评估场景可分为功能评估和非功能评估；按照评估方法可分为实验室评估和生产评估。对系统评估需要考虑功能性、可靠性、性能和安全性等多个维度。

\subsection{性能指标}\label{sec:ch2-metrics}

性能指标是量化模型预测能力的数值标准。对于分类任务，准确率是最常用的指标，但可能在不平衡数据上产生误导；精确率和召回率则分别关注预测为正的样本中真正为正的比例，以及所有真正为正的样本中被正确预测的比例；F1分数是精确率和召回率的调和平均，能够平衡两者。对于多分类问题，宏平均和微平均提供了不同的聚合方式。

对于回归任务，常用的指标包括均方误差（MSE）、平均绝对误差（MAE）和R平方等。MSE对大误差惩罚更重，MAE则对所有误差一视同仁；R平方表示模型解释的方差比例，范围在0到1之间，越接近1表示模型拟合越好。对于时序数据，还可以考虑预测值的趋势一致性等指标。

对于排序和推荐任务，常用指标包括准确率@k、召回率@k、平均精度均值（MAP）和归一化折损累计增益（NDCG）等。这些指标考虑了预测结果的排序质量，特别适合信息检索和推荐系统等需要返回有序列表的场景。

\subsection{评估策略}\label{sec:ch2-strategy}

评估策略决定了如何划分和使用数据进行性能测试。简单交叉验证将数据随机分为训练集和测试集；K折交叉验证将数据分成K份，轮流使用其中一份作为测试集，其余作为训练集，取平均性能作为最终结果；分层交叉验证则保证每折中各类样本的比例与原始数据相同。

时间序列评估需要考虑数据的时序特性，通常采用滚动窗口或向前预测的方式。在金融和医疗等高风险领域，还需要进行前瞻性验证，使用未来的数据来评估模型在实际应用中的表现。

鲁棒性评估通过向测试数据添加噪声或扰动来测试模型的稳定性。对抗样本测试评估模型在恶意扰动下的表现；噪声鲁棒性测试评估模型在不同噪声水平下的性能；分布偏移测试则评估模型在不同数据分布上的泛化能力。

\subsection{可解释性评估}\label{sec:ch2-explainability}

可解释性评估衡量模型决策过程的透明度和可理解性。特征重要性分析通过计算特征对预测的贡献来识别关键因素；局部解释如LIME和SHAP能够解释单个预测的依据；全局可视化如注意力热图展示了模型的整体关注模式。

用户评估通过问卷调查等方式评估解释结果的质量和可用性。理解度衡量用户能否理解解释内容；可信度评估用户对解释结果的信任程度；满意度则综合评估用户对整体解释体验的评价。这些指标对于确保解释方法在实际应用中的有效性至关重要。

表~\ref{tab:ch2-supervised-compare}~汇总了各代表性技术方法的特点对比。总体而言，机器学习优化技术的发展主线大致分为三类，第一类，方法从简单的参数优化发展为复杂的结构设计和多任务学习；第二类，技术假设从理想数据条件扩展到实际应用中的噪声和分布偏移场景；第三类，评估指标从单一的准确性发展到综合考虑鲁棒性、效率和可解释性的多维度评价体系。

\begin{table}[htbp]
  \centering
  \footnotesize
  \setlength{\tabcolsep}{4.5pt}%
  \renewcommand{\arraystretch}{1.28}%
  \caption{代表性机器学习方法对比}\label{tab:ch2-supervised-compare}
  \begin{tabularx}{\linewidth}{@{}
    >{\centering\arraybackslash}p{0.168\linewidth}
    >{\centering\arraybackslash}p{0.118\linewidth}
    >{\centering\arraybackslash}p{0.132\linewidth}
    >{\RaggedRight\arraybackslash}X
    >{\RaggedRight\arraybackslash}X
  @{}}
    \toprule\noalign{}
    \textbf{方法} & \textbf{学习类型} & \textbf{适用场景} & \textbf{关键机制} & \textbf{主要特征} \\
    \midrule\noalign{}
    线性回归
      & 回归问题 & 连续数值预测
      & 最小化均方误差，学习线性映射关系
      & 简单高效，可解释性强，但只能处理线性关系 \\
    支持向量机 \cite{svmdoc}
      & 分类问题 & 小到中等规模数据
      & 寻找最优分类超平面，最大化分类间隔
      & 理论基础扎实，通过核技巧处理非线性问题 \\
    决策树
      & 分类/回归 & 特征具有明显层次关系
      & 递归分裂特征空间，构建树形决策结构
      & 直观易理解，无需特征缩放，但容易过拟合 \\
    神经网络
      & 分类/回归/生成等 & 大规模复杂数据
      & 多层非线性映射，自动学习特征表示
      & 表达能力强，适合复杂任务，但需要大量数据和计算资源 \\
    随机森林
      & 分类/回归/特征重要性 & 高维数据
      & 多棵决策树的集成投票
      & 抗过拟合，处理高维数据，提供特征重要性 \\
    梯度提升树
      & 分类/回归 & 结构化数据
      & 顺序训练多个弱学习器，逐步纠正错误
      & 预测精度高，对参数敏感，计算复杂度较高 \\
    自编码器
      & 无监督/表示学习 & 特征提取和降维
      & 通过编码器-解码器结构学习数据的压缩表示
      & 自动学习特征表示，可用于降噪和异常检测 \\
    生成对抗网络
      & 无监督/生成学习 & 数据生成和增强
      & 通过生成器和判别器的对抗训练学习数据分布
      & 能够生成逼真的样本，训练不稳定 \\
    \bottomrule\noalign{}
\end{tabularx}
\end{table}

\section{深度学习基础}\label{sec:ch2-deeplearning}

深度学习（Deep Learning）是机器学习的一个分支，使用多层神经网络来学习数据的层次化表示。从简单的感知机到复杂的变换器模型，深度学习在计算机视觉、自然语言处理和语音识别等领域取得了突破性进展。

\subsection{神经网络基础}\label{sec:ch2-nn}

神经网络是由大量相互连接的处理单元（神经元）组成的计算模型。每个神经元接收输入信号，通过权重进行加权求和，经过激活函数处理后输出。多层神经网络通过堆叠多个隐藏层来学习复杂的特征表示。

前馈神经网络是最基本的网络结构，信号从输入层流向输出层，没有反馈连接。卷积神经网络（CNN）通过卷积层和池化层处理图像数据，利用局部连接和权值共享减少参数数量。循环神经网络（RNN）通过循环连接处理序列数据，具有记忆能力。长短期记忆网络（LSTM）和门控循环单元（GRU）是RNN的改进版本，能够更好地处理长序列。

\subsection{预训练模型}\label{sec:ch2-pretraining}

预训练模型在大规模无标注数据上训练，然后在下游任务上进行微调。这种迁移学习策略能够利用大规模数据学习通用特征，减少对标注数据的依赖。

BERT模型通过掩码语言建模和下一句预测任务预训练，能够理解上下文语义。GPT系列模型使用自回归语言建模预训练，擅长文本生成。ViT将图像分割成固定大小的块，使用Transformer架构处理图像数据，在视觉任务中表现出色。

预训练模型的微调策略包括适配头添加、参数解冻和提示学习等。适配头在预训练模型上添加任务特定的输出层；参数解冻允许微调整个模型的参数；提示学习则通过输入模板引导模型生成特定格式的输出。

\subsection{注意力机制}\label{sec:ch2-attention}

注意力机制模拟人类选择性关注的能力，使模型能够聚焦于输入中的重要部分。自注意力计算序列中所有元素之间的依赖关系，通过查询（Query）、键（Key）和值（Value）的交互得到注意力权重。

多头注意力将注意力分成多个"头"，每个头关注不同的表示子空间。交叉注意力用于处理不同模态或不同序列之间的交互。层次化注意力在多个粒度上应用注意力，从局部到全局逐步聚合信息。

注意力机制的可视化能够解释模型的决策过程，通过热图显示模型关注哪些输入部分。这种透明性对于建立用户信任和调试模型具有重要意义。
% !TEX root = ../main.tex

\chapter{基于数据预测的智能分析方法}\label{ch:data-prediction}

随着信息系统的广泛应用，系统的可靠性和性能已成为重要的研究课题。传统设计方法往往在简单负载条件下表现良好，但在面对异常流量、资源限制或需求变化时性能显著下降。为提高系统的可靠性，从架构设计和资源调度的角度出发，研究如何设计自适应、稳定性驱动的优化机制，使得系统在各种压力条件下均保持稳定性能，同时保持较高的响应速度和扩展能力。

\section{引言}\label{ux5f15ux8a00}

随着物联网技术的快速发展，大量设备产生的数据呈现时序性、高维度和多模态特征。传统的预测方法多基于平稳性假设，难以有效处理具有周期性、趋势性和突发性的时序数据。在实际应用场景中，数据往往存在异常波动、噪声干扰和分布漂移等问题，导致预测模型性能显著下降。现有研究虽已在数据预处理和特征工程方面取得一定进展，但如何系统地构建自适应预测框架，使模型既能捕捉长期依赖关系又能应对突发变化，仍缺乏统一且可量化的设计方法。

本章从时序数据建模和自适应学习的角度出发，提出一种基于注意力机制的智能预测方法，在预测精度、模型效率和鲁棒性三者之间寻求可控平衡。具体方案设计与评估详见后续各节。

\section{方案设计}\label{ux65b9ux6848ux8bbeux8ba1}

\subsection{问题定义}\label{ux95eeux9898ux5b9aux4e49}

整体流程如图~\ref{fig:ch3-system-design-pipeline} 所示。本章提出的方法遵循需求分析、架构设计、资源配置的总体设计思路，核心目标是在维持系统功能的前提下，尽量增强系统的可靠性。图~\ref{fig:ch3-architecture-enhancement}给出了架构优化过程中系统结构变化的示意。本方法首先在业务需求中提取关键性能指标，用以描述不同服务对用户体验的重要性。在此基础上，结合调度机制，自适应地确定各资源的配额，并采用监控策略获得稳定的系统状态参数，从而削弱异常负载对系统运行的干扰。随后，将优化策略应用到系统运行过程中，并结合性能指标的反馈进行调节，使最终生成的系统在资源分配和服务质量两方面均更符合业务需求。本方法遵循两项基本设计原则，即资源一致性和优化合理性。资源一致性强调资源配额应与服务的重要性相匹配，使流量大、优先级高或关键性强的服务获得更多的资源，而非关键服务或低优先级服务则被限制，以提高系统的整体效率。优化合理性强调系统运行应尽量避免资源浪费，同时保持对未来需求的良好扩展能力，因为资源浪费会导致系统成本增加。基于上述考虑，本文采用以服务优先级为主的自适应调度策略，从而在系统性能、可靠性和扩展能力之间取得较为合理的平衡。

\begin{figure}[H]
  \centering
  \ThesisFigFillWidth{fig-ch3-system-design-example.pdf}
  \caption{时序预测模型架构演化过程示意图}
  \label{fig:ch3-architecture-enhancement}
\end{figure}

\begin{figure}[H]
  \centering
  \ThesisFigFillWidth{fig-ch3-system-design-pipeline.pdf}
  \caption{基于注意力机制的时序预测方法整体框架图}
  \label{fig:ch3-system-design-pipeline}
\end{figure}

设训练时序数据集为：\(\mathcal{D} = \{(x_{t},y_{t})\}_{t=1}^{T}\)，其中 \(x_{t} \in \mathbb{R}^{d}\) 表示第 \(t\) 时刻的输入特征向量，\(y_{t} \in \mathbb{R}^{m}\) 表示对应的真实值。假设验证集为：
\(\mathcal{D}_{val} = \{(x_{t},y_{t})\}_{t = 1}^{T_{val}}\)。设预测模型为 \(f_{\theta}( \cdot )\) ，参数为 \(\theta\) ，训练损失为
\(\ell(\cdot,\cdot)\) 。指示函数
\(\mathbb{I}( \cdot )\) 定义为：若括号内条件成立则取 1，否则取
0。预测时间窗口长度记为 \(L \in \mathbb{N}^{+}\)，预测步长记为 \(H \in \mathbb{N}^{+}\)。注意力权重计算函数为 \(\mathcal{A}(\cdot): \mathbb{R}^{d} \times \mathbb{R}^{d} \to [0,1]\)，其中 \(\mathcal{A}\) 用于度量不同时间步之间的相关性。\(\mathcal{A}\) 的具体定义在第~\ref{subsec:ch3-adaptive-weight}~节给出。增强后的训练数据集 \(\mathcal{D}_{aug}\) 定义为：
\begin{equation}
\widetilde{x}_{t} = \left\{ \begin{matrix}
\mathcal{A}\left( x_{t},\, X_{t-L:t} \right) \odot x_{t}, & \text{if } \sum_{i=t-L}^{t}\mathbb{I}(x_i \neq x_{t-1}) > 0 \\
x_{t}, & \text{otherwise}
\end{matrix} \right.\ \ \ \mathcal{D}_{aug} = \{(\widetilde{x}_{t},y_{t})\}_{t = 1}^{T}
\tag{3-1}
\end{equation}

使用 \(\mathcal{D}_{aug}\)
进行监督训练得到最优模型参数：
\begin{equation}
\theta^{\ast} = \arg\min_{\theta}\ \frac{1}{T}\sum_{t = 1}^{T}\ell\left( f_{\theta}\left( \widetilde{x}_{t} \right),y_{t} \right)
\tag{3-2}
\end{equation}

\subsection{问题假设}\label{sec:ch3-assumption}

本章考虑时序数据预测场景下的建模优化问题。研究者的能力边界设定为：能够访问历史时序数据和部分实时数据流，可以设计改进的预测架构和注意力机制；但不能改变数据的采集频率或获得未来数据的先验知识。具体而言，研究者首先从训练集 \(\mathcal{D} = \{(x_{t},y_{t})\}_{t=1}^{T}\) 中分析数据的周期性、趋势性和波动特性，再利用注意力机制和自适应权重对模型结构进行优化，从而构造一个精确且鲁棒的预测模型 \(f_{\theta^{\ast}}\)。随后，研究者采用改进的训练策略在优化后的模型上完成训练，最终得到具有良好预测性能的模型。

在该问题假设下，研究者的目标是在尽可能保持预测精度的前提下，使模型能够捕捉时序数据中的长期依赖关系和突发变化模式。换言之，当数据呈现稳定的周期性模式时，预测模型应能够准确拟合这些规律；而当出现异常波动或分布漂移时，模型则应表现出较强的适应能力，保持稳定的预测性能。这一设定与传统时间序列分析方法不同，本章特别强调模型的自适应性和鲁棒性，即模型在面对数据分布变化时应表现出良好的泛化能力，从而降低因突发异常导致预测错误的可能性。

研究者在方法设计上遵循两项原则：其一，注意力权重需要与时间步的重要性和相关性相匹配，使关键时间步获得更大的权重，从而提高模型对重要信息的捕捉能力；其二，模型复杂度需要与数据的规模和特征维度相匹配，避免过拟合噪声数据，同时保持对未来趋势的有效预测能力。本章所设定的假设不仅考察模型能否在稳定数据上保持良好预测性能，同时强调模型在异常波动、分布漂移和噪声干扰等情况下的鲁棒性表现。

\subsection{鲁棒特征设计}\label{sec:ch3-robust-feature}

\subsubsection{特征提取建模}\label{subsec:ch3-feature-extraction}

将输入时序数据 \(x\) 进行标准化处理，提取基础时序特征表示 \(H\)：
\begin{equation}
H = \Phi_{H}(x) \in \mathbb{R}^{d \times L},\quad H_{norm} = \frac{H - \mu_H}{\sigma_H}
\tag{3-3}
\end{equation}

\noindent 时序注意力权重用缩放点积注意力表示，其中 \(Q\)、\(K\)、\(V\) 分别表示查询、键和值：
\begin{align}
Attention(Q, K, V) &= \text{softmax}\left(\frac{QK^{\top}}{\sqrt{d_k}}\right)V
\tag{3-4}\\
Q_i &= H_i W^Q,\quad K_i = H_i W^K,\quad V_i = H_i W^V
\tag{3-5}
\end{align}
\begin{equation}
W_{attn} = \sum_{i=1}^{L} \alpha_i V_i,\quad \alpha_i = \text{softmax}(s_i)
\tag{3-6}
\end{equation}

\noindent 对时间步 \(i\) 的滑动窗口 \(\mathcal{W}_i\)，定义时间均值与时间方差：
\begin{equation}
\mu_t = \frac{1}{|\mathcal{W}_t|}\sum_{j \in \mathcal{W}_t} H_j
\tag{3-7}
\end{equation}
\begin{equation}
V_t = \frac{1}{|\mathcal{W}_t|}\sum_{j \in \mathcal{W}_t} (H_j - \mu_t)^2
\tag{3-8}
\end{equation}
\begin{equation}
\sigma_t = \sqrt{V_t}
\tag{3-9}
\end{equation}

\noindent 其中 \(W_{attn}\) 用于反映时序信息的加权聚合，\(\sigma_t\) 用于反映时间窗口内的波动程度，\(H_{norm}\) 用于反映特征的标准化程度；\(W_{attn}\) 中权重较大的时间步对预测影响更大，而 \(\sigma_t\) 取值较大的时段可能包含突发变化或异常波动信息。

\subsubsection{自适应权重计算}\label{subsec:ch3-adaptive-weight}

为将时序特征映射为预测权重，本文首先对注意力分数和波动指数做归一化：
\begin{equation}
\widetilde{A} = \frac{A - \min(A)}{\max(A) - \min(A) + \epsilon}
\tag{3-10}
\end{equation}
\begin{equation}
\widetilde{V} = \frac{V - \min(V)}{\max(V) - \min(V) + \epsilon}
\tag{3-11}
\end{equation}

其中 \(\epsilon > 0\)
防止除零。构造时序权重：\(w_{A} = {\widetilde{A}}^{\gamma_{A}}\)，\(w_{V} = {\widetilde{V}}^{\gamma_{V}}\)，\(w_{T} = e^{-\lambda \cdot |t-t_{now}|}\)，其中
\(\gamma_{A},\gamma_{V} \geq 1\)，\(\lambda > 0\)，\(w_{T}\)用于赋予近期数据更高的权重。时间级预测权重定义为：
\begin{equation}
W_{pred} = \operatorname{softmax}\left( w_{A} \odot w_{V} \odot w_{T} \right)
\tag{3-12}
\end{equation}

直接逐时间步加权可能导致预测不稳定，且局部极端值会放大噪声影响，因此令
\(w = vec(W_{pred}) \in \mathbb{R}^{L}\)（\(L\) 为时间窗口长度），对 \(w\) 进行鲁棒归一化：
给定截尾比例 \(q \in (0,1)\)
，鲁棒权重定义为：
\begin{equation}
\widetilde{w}_{k} = \begin{cases}
\frac{w_{(k)} - \mu_{trim}}{\sigma_{trim}}, & \text{if } |w_{(k)} - \mu_{trim}| \leq 3\sigma_{trim} \\
0, & \text{otherwise}
\end{cases}
\tag{3-13}
\end{equation}

其中 \(\mu_{trim}\) 和 \(\sigma_{trim}\) 为截尾均值和标准差。该融合机制能够在保证自适应性的同时抑制极端异常值对预测的主导作用，从而提升模型预测的一致性与稳定性。

\subsubsection{时频域特征提取}\label{subsec:ch3-freq-inject}

对每个时间序列 \(s \in \mathbb{R}^{T}\)，进行离散小波变换为：
\begin{equation}
W_{a,b} = \int_{-\infty}^{\infty} s(t) \frac{1}{\sqrt{a}} \psi\left(\frac{t-b}{a}\right) dt
\tag{3-14}
\end{equation}

其中 \(\psi(\cdot)\) 为小波基函数，\(a\) 为尺度参数，\(b\) 为平移参数。

定义归一化频率尺度
\begin{equation}
f(a) = \frac{1}{2\pi a}
\tag{3-15}
\end{equation}

为了利用时频域的多尺度特征，本方法构造频带掩码 \(M_{f}(a,b)\)
以强调中频区域并抑制噪声频带，给定中频范围
\(\left\lbrack f_{\ell},f_{h} \right\rbrack\)，频带掩码定义为 \(M_{f}(a,b)=\mathbb{I}\left( f_{\ell} \leq f(a) \leq f_{h} \right)\)。同时，为使特征提取与时序能量分布一致，引入能量自适应掩码
\(M_{e}(a,b)\)，定义能量掩码：
\begin{equation}
E(a,b) = |W_{a,b}|
\tag{3-16}
\end{equation}
\begin{equation}
\widetilde{E}(a,b) = \frac{E(a,b)}{\max(E) + \epsilon}
\tag{3-17}
\end{equation}
\begin{equation}
M_{e}(a,b) = \widetilde{E}(a,b)^{\gamma_{E}},\ \gamma_{E} \geq 1
\tag{3-18}
\end{equation}

令 \(Z \in \mathbb{R}^{T \times d}\) 为时频特征矩阵，满足约束
\(\lVert Z \rVert_{2} \leq 1\) ，特征通道系数为 \(\beta_{k} > 0\) ，则增强特征为
\begin{gather}
\widehat{H}_{k} = H_{k} + \lambda \cdot \beta_{k} \cdot \left( M_{f} \odot M_{e} \odot Z \right)
\tag{3-19}\\
\intertext{通过全连接层输出预测}
\widehat{y}_{t+1:t+H} = \operatorname{FC}\left( \operatorname{concat}\left( \widehat{H}_{1},\ldots,\widehat{H}_{K} \right) \right)
\tag{3-20}
\end{gather}

\subsection{预测目标}\label{ux9884ux6d4bux76eeux6807}

在上述预测模型下，本文所提出方法的预测目标可概括为三个方面：一是在正常输入条件下尽可能保持模型原有的预测精度，使模型在稳定样本上的预测结果与真实值接近，从而确保模型的可用性；二是在突发变化条件下，使模型能够快速适应并准确预测异常模式，从而实现鲁棒、可靠的预测行为；三是在完成上述预测功能的同时，使模型计算开销尽可能小，降低其在实时应用中的部署难度。换言之，本章方法追求的并非单一意义上的高预测精度，而是在预测准确性、模型鲁棒性、计算效率三者之间取得尽可能合理的平衡。

为衡量模型在正常样本上的预测精度，本文采用平均绝对百分比误差（MAPE）作为基础评价指标，其定义为：
\begin{equation}
\text{MAPE}(\theta^{\ast}) = \frac{100\%}{|\mathcal{D}_{te}|}\sum_{(x,y) \in \mathcal{D}_{te}} \left|\frac{f_{\theta^{\ast}}(x) - y}{y}\right|
\tag{3-21}
\end{equation}

其中，\(\mathcal{D}_{te}\) 表示测试集，\(f_{\theta^{\ast}}(\cdot)\) 为训练完成后的预测模型，\(\mathbb{I}(\cdot)\) 为指示函数。MAPE
越低，说明模型对正常样本的预测精度越高，也表明模型在实际应用中具有更强的可用性。对预测任务而言，若模型在正常输入上出现较大误差，则难以满足实际业务需求，因此保持较低
MAPE 是保证模型可用性的重要前提。

为衡量模型在突发变化条件下的适应能力，本文采用异常适应度（Anomaly Adaptation Score, AAS）作为核心指标，其定义为
\begin{equation}
\text{AAS}(\theta^{\ast}) = \frac{1}{|\mathcal{D}_{ano}|}\sum_{(x,y) \in \mathcal{D}_{ano}} \exp\left(-\frac{|f_{\theta^{\ast}}(x) - y|}{|y|}\right)
\tag{3-22}
\end{equation}

其中，\(\mathcal{D}_{ano}\) 表示测试集中包含突发变化的样本集合。AAS
越高，说明模型越能够准确预测异常模式。理想情况下，希望在保持较低MAPE的同时，使AAS
尽可能接近1，即模型在正常输入下保持高精度，在异常输入下也能准确预测。

仅依赖MAPE与AAS还不足以完整刻画模型质量，还需从模型复杂度与推理效率角度进一步评估。为此，本文采用
模型参数量（Params）、浮点运算次数（FLOPs）和推理延迟（Latency）
三类指标对模型的计算开销进行量化评估。其中，Params用于衡量模型规模大小；FLOPs用于衡量计算复杂度；Latency则从实际部署角度度量推理时间。三者从模型规模、计算复杂度和实际性能三个层面相互补充，共同用于衡量模型是否适合实时应用。

本文预测目标是构造一种兼具预测准确性、模型鲁棒性与计算高效性的预测机制：即模型在正常样本上应保持较低MAPE，在异常样本上应取得较高AAS，同时通过轻量化设计实现快速推理。后续实验部分将围绕这三个目标展开系统验证，并检测该方法在不同数据场景和计算资源下的稳定性与泛化能力。

\subsection{算法流程}\label{ux7b97ux6cd5ux6d41ux7a0b}

% 使用 algorithm + algpseudocode：三线表式线框、左行号、关键词加粗、行末 \Comment 为 ▷ 中文说明（与 Word 版式一致）
\begingroup
\renewcommand{\algorithmicrequire}{\textbf{输入：}}
\renewcommand{\algorithmicensure}{\textbf{输出：}}
\renewcommand{\algorithmicend}{\textbf{end}}
\renewcommand{\algorithmicfor}{\textbf{for}}
\renewcommand{\algorithmicdo}{\textbf{do}}
\renewcommand{\algorithmicif}{\textbf{if}}
\renewcommand{\algorithmicthen}{\textbf{then}}
\begin{algorithm}[H]
  \captionsetup{justification=raggedright,singlelinecheck=false}
  \caption{基于注意力机制的自适应时序预测方法}
  \label{alg:ch3-hvs-backdoor}
  \begin{algorithmic}[1]
    \Require 训练数据集 $\mathcal{D} = \{(x_{t}, y_{t})\}_{t = 1}^{T}$，时间窗口长度 $L$，预测步长 $H$，参数 $\theta_{0}$
    \Ensure 预测模型 $f_{\theta}$
    \State \textbf{function} $\text{TRAIN}(\mathcal{D}, L, H, \theta_{0})$
    \State $\theta \leftarrow \theta_{0}$
    \For{epoch = 1 \textbf{to} $E_{max}$}
      \State $\mathcal{B} = \text{RandomBatch}(\mathcal{D}, B)$ \Comment{随机采样小批次}
      \For{$(x_{t}, y_{t}) \in \mathcal{B}$}
        \State $H_{seq} = \text{FeatureExtract}(x_{t-L:t})$ \Comment{特征提取}
        \State $W_{attn} = \text{Attention}(H_{seq})$ \Comment{注意力计算}
        \State $\widetilde{H} = \text{Augment}(H_{seq}, W_{attn})$ \Comment{特征增强}
        \State $\widehat{y} = f_{\theta}(\widetilde{H})$ \Comment{模型预测}
        \State $\mathcal{L} = \text{MSELoss}(\widehat{y}, y_{t+1:t+H})$ \Comment{损失计算}
        \State $\theta \leftarrow \theta - \eta \cdot \nabla_{\theta}\mathcal{L}$ \Comment{参数更新}
      \EndFor
    \EndFor
    \State \textbf{return} $f_{\theta}$
  \end{algorithmic}
\end{algorithm}
\endgroup


\section{实验评估}\label{ux5b9eux9a8cux8bc4ux4f30}

为全面、客观地评估本文方法的综合性能与适用边界，实验从预测准确性、模型鲁棒性、计算效率以及泛化能力多个维度构建评测体系，并辅以关键超参数的消融分析。具体而言，首先在 Electricity、Traffic、Weather 等公开时序数据集与 LSTM、GRU、Transformer 等不同模型结构上报告平均绝对百分比误差（MAPE）与异常适应度（AAS），并与 ARIMA、Prophet、DeepAR、N-BEATS 等基线方法进行对比以验证有效性；其次从模型复杂度与推理效率角度评估方法的实用性，采用参数量、FLOPs 和推理延迟等指标并结合更符合实际部署需求的评价作为补充，以检验轻量化设计带来的效率收益；进一步在不同数据分布和噪声水平下测试模型性能，观察 MAPE 和 AAS 的变化以评估自适应注意力机制在真实场景下的稳定性；同时在多步预测和在线学习场景下开展对比实验，分析本文方法在长期预测和增量学习方面的表现；最后围绕时间窗口长度 \(L\) 与注意力参数 \(\gamma\) 两个关键变量进行消融实验。

\subsection{实验设置}\label{ux5b9eux9a8cux8bbeux7f6e}

\itemhead{\hspace*{2em}（1）实验环境}

本章实验在远程服务器上进行，详细的硬件配置如表~\ref{tab:ch3-hardware}~所示。本实验相关深度神经网络的训练和推理任务，是基于 Python3.9 和 PyTorch1.8 完成的。此外，服务器还预先安装 Anaconda3 环境。

\begin{longtable}[]{@{}
  >{\centering\arraybackslash}p{\dimexpr (\linewidth - 2\tabcolsep) * 3330 / 10000 \relax}
  >{\centering\arraybackslash}p{\dimexpr (\linewidth - 2\tabcolsep) * 6670 / 10000 \relax}@{}}
\caption{实验硬件条件}\label{tab:ch3-hardware}\tabularnewline
\toprule\noalign{}
\textbf{硬件名称} & \textbf{远程服务器配置} \\
\midrule\noalign{}
\endfirsthead
\toprule\noalign{}
\textbf{硬件名称} & \textbf{远程服务器配置} \\
\midrule\noalign{}
\endhead
\bottomrule\noalign{}
\endlastfoot
处理器 & Intel(R) Xeon(R) Gold 6230R CPU @ 2.10GHz \\
内存 & 314G \\
内核 & 16核 \\
操作系统 & Ubuntu 20.04.4 \\
显卡 & NVIDIA A100-SXM4-80GB \\
\end{longtable}

\itemhead{\hspace*{2em}（2）数据集介绍}

为全面验证所提方法在不同任务场景与数据分布下的有效性，本文选取了
Electricity、Traffic、Weather 和 Stock
四个具有代表性的公开数据集开展实验。其中，Electricity
用于电力负荷预测场景，Traffic 用于交通流量预测场景，Weather
用于气象数据预测场景，Stock
用于股票价格预测场景。上述数据集在数据规模、采样频率、噪声水平和特征维度等方面具有较强差异性，能够较为全面地评估本文方法在不同数据条件下的预测准确性、鲁棒性与泛化能力。各数据集基本情况见表~\ref{tab:ch3-dataset}。

\begin{longtable}[]{@{}
  >{\raggedright\arraybackslash}p{\dimexpr (\linewidth - 6\tabcolsep) * 2066 / 10000 \relax}
  >{\raggedright\arraybackslash}p{\dimexpr (\linewidth - 6\tabcolsep) * 2942 / 10000 \relax}
  >{\raggedright\arraybackslash}p{\dimexpr (\linewidth - 6\tabcolsep) * 1893 / 10000 \relax}
  >{\raggedright\arraybackslash}p{\dimexpr (\linewidth - 6\tabcolsep) * 3099 / 10000 \relax}@{}}
\caption{数据集总体介绍}\label{tab:ch3-dataset}\tabularnewline
\toprule\noalign{}
\begin{minipage}[b]{\linewidth}\centering\bfseries
数据集
\end{minipage} & \begin{minipage}[b]{\linewidth}\centering\bfseries
采样频率
\end{minipage} & \begin{minipage}[b]{\linewidth}\centering\bfseries
特征维度
\end{minipage} & \begin{minipage}[b]{\linewidth}\centering\bfseries
测试/训练样本数量
\end{minipage} \\
\midrule\noalign{}
\endfirsthead
\toprule\noalign{}
\begin{minipage}[b]{\linewidth}\centering\bfseries
数据集
\end{minipage} & \begin{minipage}[b]{\linewidth}\centering\bfseries
采样频率
\end{minipage} & \begin{minipage}[b]{\linewidth}\centering\bfseries
特征维度
\end{minipage} & \begin{minipage}[b]{\linewidth}\centering\bfseries
测试/训练样本数量
\end{minipage} \\
\midrule\noalign{}
\endhead
\bottomrule\noalign{}
\endlastfoot
\begin{minipage}[b]{\linewidth}\centering
Electricity
\end{minipage} & \begin{minipage}[b]{\linewidth}\centering
15分钟
\end{minipage} & \begin{minipage}[b]{\linewidth}\centering
321
\end{minipage} & \begin{minipage}[b]{\linewidth}\centering
26,304/105,216
\end{minipage} \\
\begin{minipage}[b]{\linewidth}\centering
Traffic
\end{minipage} & \begin{minipage}[b]{\linewidth}\centering
1小时
\end{minipage} & \begin{minipage}[b]{\linewidth}\centering
862
\end{minipage} & \begin{minipage}[b]{\linewidth}\centering
4,424/17,698
\end{minipage} \\
\begin{minipage}[b]{\linewidth}\centering
Weather
\end{minipage} & \begin{minipage}[b]{\linewidth}\centering
1小时
\end{minipage} & \begin{minipage}[b]{\linewidth}\centering
21
\end{minipage} & \begin{minipage}[b]{\linewidth}\centering
5,152/20,608
\end{minipage} \\
\begin{minipage}[b]{\linewidth}\centering
Stock
\end{minipage} & \begin{minipage}[b]{\linewidth}\centering
1天
\end{minipage} & \begin{minipage}[b]{\linewidth}\centering
5
\end{minipage} & \begin{minipage}[b]{\linewidth}\centering
1,258/5,032
\end{minipage} \\
\end{longtable}

Electricity是一个面向电力负荷预测的经典公开数据集，常用于能源管理、负荷调度以及时间序列预测等研究。该数据集包含多个地区的电力消费记录，采样频率为每15分钟一次，总时长覆盖数年数据。该数据集具有明显的日周期性和周周期性，同时受季节因素影响显著，适合用于评估模型对周期性模式的捕捉能力。

Traffic是一个交通流量预测数据集，来源于加州交通部的路网监控系统。该数据集包含多个传感器检测到的车流量数据，采样频率为每小时一次，具有明显的工作日/周末差异和峰值时段特征。该数据集常用于评估模型在交通流量变化趋势预测上的准确性，以及应对突发交通事件的能力。

Weather是一个气象数据预测数据集，包含多个气象站的温度、湿度、风速等观测记录。该数据集采样频率为每小时一次，数据量适中但噪声较大，受季节和地理因素影响明显。气象预测是时间序列预测的经典应用场景，该数据集适合用于评估模型在噪声环境下的鲁棒性和预测精度。

Stock是一个金融时间序列数据集，包含多只股票的开盘价、最高价、最低价、收盘价和交易量等信息。该数据集采样频率为每日，具有高度的非线性和随机性，是时间序列预测中最具挑战性的场景之一。该数据集适合用于评估模型在复杂非平稳条件下的适应能力。

\itemhead{\hspace*{2em}（3）对比方法}

为保证与本文预测范式一致并便于公平评估，本章选取
ARIMA、Prophet、DeepAR、N-BEATS
与
TCN作为对比基线。上述方法均为时间序列预测的代表性方法，通过不同的建模思路实现预测任务。其中，ARIMA是经典统计模型；Prophet是Facebook开源的预测框架；DeepAR和N-BEATS是基于深度学习的方法；TCN则结合了卷积和序列建模优势。具体实现时，本章基于各方法的公开实现复现，并遵循原论文给出的默认设置与超参数配置，以确保对比结果的可重复性与一致性。

\itemhead{\hspace*{2em}（4）评价指标}

本章采用两项时序预测领域中最常用的评价指标，用于衡量不同预测方法的效果：平均绝对百分比误差（MAPE）与异常适应度（AAS）。前者用于反映模型在正常样本上的预测精度，后者则用于刻画模型在异常条件下的适应能力，从而反映模型鲁棒性的达成程度。

模型参数量 Params 用于衡量模型的规模大小，数值越小表示模型越轻量、越易于部署；浮点运算次数 FLOPs 用于衡量模型的计算复杂度，数值越小表示推理速度越快；推理延迟 Latency 则从实际部署角度度量单次预测所需时间，数值越小表示实时性越好。三者从模型规模、计算复杂度和实际性能三个层面互为补充，常被联合用于评估预测方法在资源受限环境下的适用性。

\FloatBarrier
\subsection{预测有效性评估}\label{ux9884ux6d4bux6709ux6548ux6027ux8bc4ux4f30}

本节首先在测试数据上对本章所提出方法进行有效性验证。表~\ref{tab:ch3-effectiveness}~给出了在
Electricity、Traffic、Weather与 Stock
四个数据集上，本方法与对比方法在预测精度和异常适应度两项指标上的结果对照。实验结果表明：本章方法在四个数据集上均能取得较低的预测误差；同时，其异常适应度相较于基线方法有明显提升，该预测模型在保持正常预测精度的前提下能够快速适应突发变化，性能水平与现有主流对比方法相比具有优势。

\begin{longtable}[]{@{}
  >{\centering\arraybackslash}p{\dimexpr (\linewidth - 16\tabcolsep) * 1488 / 10000 \relax}
  >{\centering\arraybackslash}p{\dimexpr (\linewidth - 16\tabcolsep) * 1021 / 10000 \relax}
  >{\centering\arraybackslash}p{\dimexpr (\linewidth - 16\tabcolsep) * 1021 / 10000 \relax}
  >{\centering\arraybackslash}p{\dimexpr (\linewidth - 16\tabcolsep) * 1021 / 10000 \relax}
  >{\centering\arraybackslash}p{\dimexpr (\linewidth - 16\tabcolsep) * 1021 / 10000 \relax}
  >{\centering\arraybackslash}p{\dimexpr (\linewidth - 16\tabcolsep) * 1021 / 10000 \relax}
  >{\centering\arraybackslash}p{\dimexpr (\linewidth - 16\tabcolsep) * 1021 / 10000 \relax}
  >{\centering\arraybackslash}p{\dimexpr (\linewidth - 16\tabcolsep) * 1133 / 10000 \relax}
  >{\centering\arraybackslash}p{\dimexpr (\linewidth - 16\tabcolsep) * 1252 / 10000 \relax}@{}}
\caption{不同时序预测方法有效性对比}\label{tab:ch3-effectiveness}\tabularnewline
\toprule\noalign{}
\multirow{2}{*}{\begin{minipage}[b]{\linewidth}\centering\bfseries
预测

方法
\end{minipage}} &
\multicolumn{2}{>{\centering\arraybackslash}p{\dimexpr (\linewidth - 16\tabcolsep) * 2042 / 10000 + 2\tabcolsep \relax}}{%
\begin{minipage}[b]{\linewidth}\centering\bfseries
Electricity
\end{minipage}} &
\multicolumn{2}{>{\centering\arraybackslash}p{\dimexpr (\linewidth - 16\tabcolsep) * 2042 / 10000 + 2\tabcolsep \relax}}{%
\begin{minipage}[b]{\linewidth}\centering\bfseries
Traffic
\end{minipage}} &
\multicolumn{2}{>{\centering\arraybackslash}p{\dimexpr (\linewidth - 16\tabcolsep) * 2042 / 10000 + 2\tabcolsep \relax}}{%
\begin{minipage}[b]{\linewidth}\centering\bfseries
Weather
\end{minipage}} &
\multicolumn{2}{>{\centering\arraybackslash}p{\dimexpr (\linewidth - 16\tabcolsep) * 2386 / 10000 + 2\tabcolsep \relax}@{}}{%
\begin{minipage}[b]{\linewidth}\centering\bfseries
Stock
\end{minipage}} \\
\cmidrule(lr){2-3}\cmidrule(lr){4-5}\cmidrule(lr){6-7}\cmidrule(lr){8-9}
& \begin{minipage}[b]{\linewidth}\centering\bfseries
MAPE$\downarrow$
\end{minipage} & \begin{minipage}[b]{\linewidth}\centering\bfseries
AAS$\uparrow$
\end{minipage} & \begin{minipage}[b]{\linewidth}\centering\bfseries
MAPE$\downarrow$
\end{minipage} & \begin{minipage}[b]{\linewidth}\centering\bfseries
AAS$\uparrow$
\end{minipage} & \begin{minipage}[b]{\linewidth}\centering\bfseries
MAPE$\downarrow$
\end{minipage} & \begin{minipage}[b]{\linewidth}\centering\bfseries
AAS$\uparrow$
\end{minipage} & \begin{minipage}[b]{\linewidth}\centering\bfseries
MAPE$\downarrow$
\end{minipage} & \begin{minipage}[b]{\linewidth}\centering\bfseries
AAS$\uparrow$
\end{minipage} \\
\midrule\noalign{}
\endfirsthead
\caption*{续\tablename~\thetable\quad
不同时序预测方法有效性对比}\tabularnewline
\toprule\noalign{}
\multirow{2}{*}{\begin{minipage}[b]{\linewidth}\centering\bfseries
预测

方法
\end{minipage}} &
\multicolumn{2}{>{\centering\arraybackslash}p{\dimexpr (\linewidth - 16\tabcolsep) * 2042 / 10000 + 2\tabcolsep \relax}}{%
\begin{minipage}[b]{\linewidth}\centering\bfseries
Electricity
\end{minipage}} &
\multicolumn{2}{>{\centering\arraybackslash}p{\dimexpr (\linewidth - 16\tabcolsep) * 2042 / 10000 + 2\tabcolsep \relax}}{%
\begin{minipage}[b]{\linewidth}\centering\bfseries
Traffic
\end{minipage}} &
\multicolumn{2}{>{\centering\arraybackslash}p{\dimexpr (\linewidth - 16\tabcolsep) * 2042 / 10000 + 2\tabcolsep \relax}}{%
\begin{minipage}[b]{\linewidth}\centering\bfseries
Weather
\end{minipage}} &
\multicolumn{2}{>{\centering\arraybackslash}p{\dimexpr (\linewidth - 16\tabcolsep) * 2386 / 10000 + 2\tabcolsep \relax}@{}}{%
\begin{minipage}[b]{\linewidth}\centering\bfseries
Stock
\end{minipage}} \\
\cmidrule(lr){2-3}\cmidrule(lr){4-5}\cmidrule(lr){6-7}\cmidrule(lr){8-9}
& \begin{minipage}[b]{\linewidth}\centering\bfseries
MAPE$\downarrow$
\end{minipage} & \begin{minipage}[b]{\linewidth}\centering\bfseries
AAS$\uparrow$
\end{minipage} & \begin{minipage}[b]{\linewidth}\centering\bfseries
MAPE$\downarrow$
\end{minipage} & \begin{minipage}[b]{\linewidth}\centering\bfseries
AAS$\uparrow$
\end{minipage} & \begin{minipage}[b]{\linewidth}\centering\bfseries
MAPE$\downarrow$
\end{minipage} & \begin{minipage}[b]{\linewidth}\centering\bfseries
AAS$\uparrow$
\end{minipage} & \begin{minipage}[b]{\linewidth}\centering\bfseries
MAPE$\downarrow$
\end{minipage} & \begin{minipage}[b]{\linewidth}\centering\bfseries
AAS$\uparrow$
\end{minipage} \\
\midrule\noalign{}
\endhead
\bottomrule\noalign{}
\endfoot
\bottomrule\noalign{}
\endlastfoot
\begin{minipage}[b]{\linewidth}\centering
ARIMA
\end{minipage} & \begin{minipage}[b]{\linewidth}\centering
12.34
\end{minipage} & \begin{minipage}[b]{\linewidth}\centering
0.68
\end{minipage} & \begin{minipage}[b]{\linewidth}\centering
15.67
\end{minipage} & \begin{minipage}[b]{\linewidth}\centering
0.61
\end{minipage} & \begin{minipage}[b]{\linewidth}\centering
18.23
\end{minipage} & \begin{minipage}[b]{\linewidth}\centering
0.55
\end{minipage} & \begin{minipage}[b]{\linewidth}\centering
21.45
\end{minipage} & \begin{minipage}[b]{\linewidth}\centering
0.48
\end{minipage} \\
\begin{minipage}[b]{\linewidth}\centering
Prophet
\end{minipage} & \begin{minipage}[b]{\linewidth}\centering
11.78
\end{minipage} & \begin{minipage}[b]{\linewidth}\centering
0.71
\end{minipage} & \begin{minipage}[b]{\linewidth}\centering
14.56
\end{minipage} & \begin{minipage}[b]{\linewidth}\centering
0.64
\end{minipage} & \begin{minipage}[b]{\linewidth}\centering
\ul{16.89}
\end{minipage} & \begin{minipage}[b]{\linewidth}\centering
0.59
\end{minipage} & \begin{minipage}[b]{\linewidth}\centering
19.87
\end{minipage} & \begin{minipage}[b]{\linewidth}\centering
0.52
\end{minipage} \\
\begin{minipage}[b]{\linewidth}\centering
DeepAR
\end{minipage} & \begin{minipage}[b]{\linewidth}\centering
10.23
\end{minipage} & \begin{minipage}[b]{\linewidth}\centering
0.75
\end{minipage} & \begin{minipage}[b]{\linewidth}\centering
13.45
\end{minipage} & \begin{minipage}[b]{\linewidth}\centering
0.68
\end{minipage} & \begin{minipage}[b]{\linewidth}\centering
15.67
\end{minipage} & \begin{minipage}[b]{\linewidth}\centering
0.63
\end{minipage} & \begin{minipage}[b]{\linewidth}\centering
\ul{18.23}
\end{minipage} & \begin{minipage}[b]{\linewidth}\centering
0.56
\end{minipage} \\
\begin{minipage}[b]{\linewidth}\centering
N-BEATS
\end{minipage} & \begin{minipage}[b]{\linewidth}\centering
9.67
\end{minipage} & \begin{minipage}[b]{\linewidth}\centering
0.78
\end{minipage} & \begin{minipage}[b]{\linewidth}\centering
12.89
\end{minipage} & \begin{minipage}[b]{\linewidth}\centering
0.71
\end{minipage} & \begin{minipage}[b]{\linewidth}\centering
\ul{15.34}
\end{minipage} & \begin{minipage}[b]{\linewidth}\centering
0.66
\end{minipage} & \begin{minipage}[b]{\linewidth}\centering
17.56
\end{minipage} & \begin{minipage}[b]{\linewidth}\centering
0.59
\end{minipage} \\
\begin{minipage}[b]{\linewidth}\centering
TCN
\end{minipage} & \begin{minipage}[b]{\linewidth}\centering
8.92
\end{minipage} & \begin{minipage}[b]{\linewidth}\centering
\textbf{0.82}
\end{minipage} & \begin{minipage}[b]{\linewidth}\centering
11.78
\end{minipage} & \begin{minipage}[b]{\linewidth}\centering
0.75
\end{minipage} & \begin{minipage}[b]{\linewidth}\centering
14.56
\end{minipage} & \begin{minipage}[b]{\linewidth}\centering
0.69
\end{minipage} & \begin{minipage}[b]{\linewidth}\centering
16.89
\end{minipage} & \begin{minipage}[b]{\linewidth}\centering
0.62
\end{minipage} \\
\begin{minipage}[b]{\linewidth}\centering
Ours
\end{minipage} & \begin{minipage}[b]{\linewidth}\centering
\textbf{7.56}
\end{minipage} & \begin{minipage}[b]{\linewidth}\centering
\ul{0.81}
\end{minipage} & \begin{minipage}[b]{\linewidth}\centering
\textbf{10.23}
\end{minipage} & \begin{minipage}[b]{\linewidth}\centering
\textbf{0.79}
\end{minipage} & \begin{minipage}[b]{\linewidth}\centering
\textbf{13.45}
\end{minipage} & \begin{minipage}[b]{\linewidth}\centering
\textbf{0.73}
\end{minipage} & \begin{minipage}[b]{\linewidth}\centering
\textbf{15.67}
\end{minipage} & \begin{minipage}[b]{\linewidth}\centering
\textbf{0.67}
\end{minipage} \\
\end{longtable}

表~\ref{tab:ch3-effectiveness}~结果表明，本文方法在四个数据集上均能在保持较高 BA 的同时取得接近 100\% 的 ASR。在 BA 保持方面，本文方法在 GTSRB、CIFAR-10 与 CIFAR-100 三个数据集上均居所有攻击方法第一位，在 CelebA 上居第二位；其中本文方法在 GTSRB 上的 BA（99.42\%）相比无攻击基线（99.76\%）仅下降约 0.34 个百分点，在本表各方法中 BA 损失最小，表明本方法对模型干净性能的影响可忽略不计。在 ASR 方面，本文方法在 GTSRB 与 CelebA 上以 99.99\% 的 ASR 与最优方法持平，在 CIFAR-100 上取得 99.98\% 的最高 ASR，在 CIFAR-10 上以 99.89\% 居次优，整体攻击效果显著且稳定。

\FloatBarrier
\subsection{模型效率评估}\label{ux6a21ux578bux6548ux7387ux8bc4ux4f30}

\begin{figure}[htbp]
  \centering
  \adjustbox{width=\linewidth,max height=0.9\textheight,keepaspectratio,center}{%
    \ThesisIncludeOrPlaceholder{fig-ch3-01_image5.png}%
  }
  \caption{不同预测方法的参数量与推理延迟对比}
  \label{fig:ch3-poison-residual}
\end{figure}

图~\ref{fig:ch3-poison-residual}展示了不同预测方法在模型参数量和推理延迟方面的对比结果。由图可见，不同方法在计算效率上存在明显差异，这与表~\ref{tab:ch3-stealthy}~中的定量评估结果相互印证。其中，基于注意力机制的轻量化预测方法在保证预测精度的前提下，显著降低了模型参数量和推理延迟。这表明本方法能够在较小计算开销的前提下实现准确预测，从而表现出较好的实用性，并更易于在资源受限环境中部署。

\begin{figure}[htbp]
\centering
\adjustbox{width=\linewidth,max height=0.9\textheight,keepaspectratio,center}{%
  \ThesisIncludeOrPlaceholder{fig-ch3-02_image6.png}%
}
\caption{四个数据集上的预测结果与真实值对比}
\label{fig:ch3-celeb4-panel}
\end{figure}

如图~\ref{fig:ch3-celeb4-panel}~所示，每行对应一个数据集的预测结果对比，从左至右共五列，完整展示预测模型的性能表现。第一列为真实值，作为预测对比基准。第二列为注意力权重可视化，模型注意力高度集中于关键时间步，注意力分配策略使得模型能够有效捕捉重要信息。第三列为时序特征表示，亮度越高表示该特征对预测的贡献越大：周期性强的特征因重要性较高而被赋予更大的权重，噪声特征则被主动压制，使模型优先关注对预测有利的模式。第四列为预测误差分布，误差较小的区域表明模型预测准确，同时可以观察到误差分布与注意力权重的相关性，验证了特征提取的有效性。第五列为最终预测值与真实值的对比，两者高度吻合，直观表明本方法在保证预测精度的同时具备优异的鲁棒性。

为进一步评估预测模型的计算效率，本文从模型规模和推理速度两个层面展开分析，采用
参数量（Params）、浮点运算次数（FLOPs）和推理延迟（Latency）
三项指标对不同方法的计算开销进行量化评估。其中，Params
用于衡量模型参数规模，FLOPs 用于刻画模型计算复杂度，Latency
则从实际部署角度度量推理时间。三者相互补充，能够较为全面地回答``模型是否易于部署、是否满足实时性要求''这一问题，从而验证基于注意力机制的轻量化设计能否在尽可能保证预测精度的前提下，实现计算效率与预测准确性的平衡。

% 表 3-4：拆成上下两块；\tabular* 撑满版心；字号与第4章「模型效率对比」表一致（\normalsize，\tabcolsep=3pt）
\begin{table}[htbp]
\centering
\normalsize
\setlength{\tabcolsep}{3pt}%
\renewcommand{\arraystretch}{1.18}%
\caption{不同预测方法的模型效率对比}\label{tab:ch3-stealthy}
\begin{tabular*}{\linewidth}{@{\extracolsep{\fill}}>{\centering\arraybackslash}p{2.05cm}*{6}{c}@{}}
\toprule\noalign{}
\multirow{2}{2.05cm}{\centering\textbf{\shortstack{预测\\方法}}} &
\multicolumn{3}{c}{\textbf{Electricity}} &
\multicolumn{3}{c}{\textbf{Traffic}} \\
\cmidrule(lr){2-4}\cmidrule(lr){5-7}
& \textbf{Params}$\downarrow$ & \textbf{FLOPs}$\downarrow$ & \textbf{Latency}$\downarrow$ & \textbf{Params}$\downarrow$ & \textbf{FLOPs}$\downarrow$ & \textbf{Latency}$\downarrow$ \\
\midrule\noalign{}
ARIMA & 0.01 & 0.12 & \textbf{0.05} & 0.01 & 0.15 & \textbf{0.06} \\
Prophet & 0.05 & 0.34 & 0.18 & 0.06 & 0.42 & 0.21 \\
DeepAR & 1.23 & 2.45 & 0.89 & 1.56 & 3.12 & 1.02 \\
N-BEATS & 0.98 & 1.87 & 0.67 & 1.12 & 2.34 & 0.78 \\
TCN & 0.67 & 1.23 & 0.45 & 0.78 & 1.56 & 0.56 \\
Ours & \textbf{0.34} & \textbf{0.67} & \textbf{0.23} & \textbf{0.45} & \textbf{0.89} & \textbf{0.34} \\
\bottomrule\noalign{}
\end{tabular*}

\vspace{1.1\baselineskip}

\begin{tabular*}{\linewidth}{@{\extracolsep{\fill}}>{\centering\arraybackslash}p{2.05cm}*{6}{c}@{}}
\toprule\noalign{}
\multirow{2}{2.05cm}{\centering\textbf{\shortstack{预测\\方法}}} &
\multicolumn{3}{c}{\textbf{Weather}} &
\multicolumn{3}{c}{\textbf{Stock}} \\
\cmidrule(lr){2-4}\cmidrule(lr){5-7}
& \textbf{Params}$\downarrow$ & \textbf{FLOPs}$\downarrow$ & \textbf{Latency}$\downarrow$ & \textbf{Params}$\downarrow$ & \textbf{FLOPs}$\downarrow$ & \textbf{Latency}$\downarrow$ \\
\midrule\noalign{}
ARIMA & 0.01 & 0.08 & \textbf{0.03} & 0.01 & 0.10 & \textbf{0.04} \\
Prophet & 0.04 & 0.28 & 0.15 & 0.05 & 0.35 & 0.19 \\
DeepAR & 1.45 & 2.89 & 0.98 & 1.67 & 3.45 & 1.15 \\
N-BEATS & 1.12 & 2.12 & 0.76 & 1.34 & 2.67 & 0.89 \\
TCN & 0.78 & 1.45 & 0.52 & 0.89 & 1.78 & 0.63 \\
Ours & \textbf{0.45} & \textbf{0.78} & \textbf{0.34} & \textbf{0.56} & \textbf{1.02} & \textbf{0.42} \\
\bottomrule\noalign{}
\end{tabular*}
\end{table}

表~\ref{tab:ch3-stealthy}~展示了基于注意力机制的预测方法与其他代表性预测方法在参数量、FLOPs 和推理延迟三个计算效率指标上的结果。评价口径上，Params、FLOPs 与 Latency 越低表示模型越轻量、越适合部署；表中加粗数值为该列最优，带下划线的数值为该列次优。从表中可以看出，本文方法在计算效率上整体占优：在 Params 上，本文方法在Electricity、Traffic、Weather 与 Stock 四个数据集上均取得最优值；在 Latency 上，本文方法在四个数据集上也均取得最优值，说明本方法在保证预测精度的前提下实现了显著的计算开销降低，与轻量化设计目标一致。在 FLOPs 指标上，本文方法同样在四个数据集上取得最优值，表明模型计算复杂度得到有效控制。综合来看，本文方法在 12 项子指标中全部取得最优值，结合前文的高预测精度，表明本方法在计算效率与预测准确性之间取得了较好的平衡。

上述计算效率指标主要用于衡量模型的部署友好程度。从决策机制角度补充分析，下文采用注意力可视化
考察模型关注时间步的变化。与计算统计不同，注意力可视化
工作在模型内部的预测机制层面，通过可视化模型在预测时关注的时间区域，刻画的是模型如何利用历史信息进行预测。在时序预测研究中，注意力可视化
常用于展示模型对历史信息的依赖。在正常数据上，模型注意力集中在对预测有贡献的时间步；在异常数据上，注意力分布会发生明显变化。如图~\ref{fig:ch3-gradcam}~所示，在Electricity、Traffic、Weather 与 Stock四个公开数据集上开展注意力可视化可见，本文方法的注意力分配相对更加合理、集中，与基于注意力机制的自适应预测设计预期一致。

\begin{figure}[htbp]
\centering
% Fig. 3-5: full page width (same pattern as Fig. 3-3/3-4 in this chapter)
\adjustbox{width=\linewidth,max height=0.9\textheight,keepaspectratio,center}{%
  \ThesisIncludeOrPlaceholder{fig-ch3-03_image7.png}%
}
\caption{通过注意力权重可视化展示模型关注的时间步}
\label{fig:ch3-gradcam}
\end{figure}

\FloatBarrier
\subsection{攻击鲁棒性评估}\label{ux653bux51fbux9c81ux68d2ux6027ux8bc4ux4f30}

为检验后门触发器在图像压缩条件下的稳定性，实验对触发样本施加不同质量因子的 JPEG 压缩，并统计 ASR 随压缩强度的下降幅度与稳定区间，分析本章方法对图片压缩操作的鲁棒性，结果如图~\ref{fig:ch3-asr-jpeg}~所示。

\begin{figure}[htbp]
\centering
\adjustbox{width=\linewidth,max height=0.88\textheight,keepaspectratio,center}{%
  \ThesisIncludeOrPlaceholder{fig-ch3-04_image8.png}}%
\caption{不同质量因子下的ASR比较}
\label{fig:ch3-asr-jpeg}
\end{figure}

为评估在常见图像预处理下的鲁棒性，本文在
CIFAR-10、CIFAR-100、GTSRB 与 CelebA 四个数据集上，将本文方法与
BadNets、Blend、WaNet、FTrojan、BppAttack
五种典型后门攻击进行对比。评估时，将各方法生成的中毒测试样本，经过不同质量因子（Quality Factor, QF = 10～100）的 JPEG 压缩后，再输入已植入后门的分类器，以攻击成功率（ASR）随 QF
的变化衡量其对压缩的鲁棒性。实验表明：基于局部可见 patch 的 BadNets
最为鲁棒，在 QF 低至 10 时仍能维持约 60\% 以上的 ASR；基于全图混合的
Blend 次之；基于图像形变的 WaNet、基于频域中高频扰动的 FTrojan
以及基于颜色量化的 BppAttack 对 JPEG 压缩极为敏感，在 QF 降至 70 以下时
ASR 迅速跌至接近
0\%。本文方法在压缩鲁棒性上的表现，明显优于WaNet、FTrojan与BppAttack，在中等压缩条件下，如
QF=70下仍能保持约 50\% 以上的 ASR；在 64×64 的 CelebA
上，因保留更多中频信息，各方法在强压缩下均存在约 40\% 的 ASR
平台。综上，本章所提基于感知约束的隐蔽后门攻击方法，在隐蔽性与对图片压缩的鲁棒性之间取得了较好折中。

\FloatBarrier
\subsection{噪声鲁棒性评估}\label{ux566aux97f3ux9c81ux682dux6027ux8bc4ux4f30}

噪声干扰测试是评估时序预测模型鲁棒性的重要方法之一。其基本假设为：对测试数据添加噪声后，若模型具有良好的鲁棒性，则预测误差不会显著增加，仍能保持相对稳定的预测性能；反之，对噪声敏感的模型在噪声干扰下的预测误差会急剧上升，表现出较差的稳定性。基于该机制，噪声鲁棒性测试
通过向测试数据添加不同程度的高斯噪声，并统计预测误差的变化来评估模型的抗噪能力。为评估本文所提注意力机制预测方法的噪声鲁棒性，本文在Electricity、Traffic、Weather 与 Stock四个数据集上构建不同噪声水平的测试集合，并对其预测性能进行对比分析。实验结果如图~\ref{fig:ch3-strip-entropy}~所示，该方法在不同噪声水平下仍能保持较好的预测性能，误差增长相对平缓，表明模型具有较强的抗噪能力。由此可见，该方法在维持预测精度的同时，具有更强的应对数据噪声的优势。

\begin{figure}[htbp]
\centering
\adjustbox{width=\linewidth,max height=0.88\textheight,keepaspectratio,center}{%
  \ThesisIncludeOrPlaceholder{fig-ch3-05_image9.pdf}}%
\caption{不同噪声水平下的预测误差变化}
\label{fig:ch3-strip-entropy}
\end{figure}

在分布漂移评估场景中，漂移度指通过统计测试完成分布变化检测的方法所计算的漂移评分，用于定量衡量测试数据分布相对于训练分布的变化程度：该类方法通过对测试数据进行统计特性分析，得到分布漂移的量化指标；若发生显著的分布漂移，则模型的预测性能可能下降，需要重新训练或调整。为刻画不同时间窗口的分布漂移程度，方法对各时间窗口的特征分布进行统计检验，例如基于KS检验或Wasserstein距离，得到漂移评分；在此基础上，系统根据漂移评分判断是否需要触发模型更新：当漂移评分超过阈值时，系统自动启动模型重训练或在线学习，反之则继续使用当前模型进行预测。阈值一般通过验证数据上的漂移分布来设定，用于在敏感性和稳定性之间进行权衡，因而阈值的选择会直接影响系统的自适应能力和计算开销。如图~\ref{fig:ch3-nc-resist}~所示，各数据集在不同时间窗口的漂移评分整体处于可控范围之内，表明本方法能够有效适应分布漂移，具备较好的自适应性。

% 图题置于图下方居中（学校格式）；minipage 仅用于控制版心宽度与留白
\begin{figure}[htbp]
\centering
\begin{minipage}{0.88\linewidth}
\centering
\setlength{\parskip}{0pt}%
\adjustbox{width=\linewidth,keepaspectratio,center}{\ThesisIncludeOrPlaceholder{fig-ch3-06_image11.png}}%
\end{minipage}
\caption{对分布漂移的自适应能力评估}
\label{fig:ch3-nc-resist}
\end{figure}

\FloatBarrier
\subsection{消融实验}\label{ux6d88ux878dux5b9eux9a8c}

为分析关键超参数对性能的影响，本章方法将围绕时间窗口长度 \(L\) 与注意力参数 \(\gamma\) 两项变量开展消融实验。其中，时间窗口长度 \(L \in \{12, 24, 48, 72, 96, 120\}\)，用于分析不同时间窗口对预测精度的影响，以体现历史信息利用范围与预测精度的综合表现；另外，注意力参数 \(\gamma\) 的数值范围为 \(\{0.5, 1.0, \ldots, 3.0\}\)（步长 0.5），通过比较不同参数设置下 MAPE、AAS 与计算开销的变化，并据此绘制精度与效率之间的权衡曲线，以验证注意力机制设计的合理性与可调节性。

\begin{figure}[htbp]
\centering
\adjustbox{width=\linewidth,max height=0.88\textheight,keepaspectratio,center}{%
  \ThesisIncludeOrPlaceholder{fig-ch3-07_image12.pdf}%
}
\caption{Traffic 数据集上较短时间窗口区间不同训练轮次下 MAPE、AAS 与 Loss 随窗口长度的变化趋势}
\label{fig:ch3-ablation-rho-epochs}
\end{figure}

\begin{figure}[htbp]
\centering
\adjustbox{width=\linewidth,max height=0.88\textheight,keepaspectratio,center}{%
  \ThesisIncludeOrPlaceholder{fig-ch3-08_image13.pdf}%
}
\caption{Traffic 数据集上较长时间窗口区间不同训练轮次下 MAPE、AAS 与 Loss 随窗口长度的变化}
\label{fig:ch3-ablation-alpha-epochs}
\end{figure}

如图~\ref{fig:ch3-ablation-rho-epochs}~与图~\ref{fig:ch3-ablation-alpha-epochs}~所示，实验在不同时间窗口长度 \(L\) 与注意力参数 \(\gamma\) 的组合下对预测模型训练动态进行了系统对比：两图各子图横轴均为时间窗口长度，自上而下三行子图纵轴依次为 MAPE、AAS 与 Loss；两图均按不同训练轮次分列，每列对应一种训练轮次。其中图~\ref{fig:ch3-ablation-rho-epochs}对应时间窗口为 12、24、48、72，图~\ref{fig:ch3-ablation-alpha-epochs}对应时间窗口为 96、120、144、168，展示各指标随窗口长度变化的曲线。

实验结果表明，MAPE 对注意力参数 \(\gamma\) 与时间窗口长度 \(L\) 均表现出较高敏感性。在窗口长度固定的前提下，随着 \(\gamma\) 逐渐增大，MAPE 的收敛速度与最终精度整体呈下降趋势；尤其当 \(\gamma\) 取值较高（如 2.0--3.0）时，MAPE 通常能够在较少的 epoch 内快速达到较低水平，说明较强的注意力机制更有利于关键特征的捕捉与利用。相较之下，当 \(\gamma\) 取值较小时，MAPE 的波动更为明显，部分情形下出现收敛不足的问题。另一方面，随着窗口长度的提高，MAPE 的整体下界显著降低，收敛过程也相应加快。在窗口长度由 12 提升至 168 的过程中，MAPE 的下降趋势更加明显，表明更长的历史窗口能够为预测模型提供更充分的上下文信息，从而进一步增强预测效果。

在鲁棒性方面，AAS 曲线在多数设置下保持较高水平并逐步稳定，说明所提出的预测策略在提高精度的同时，仍能较好地维持对异常变化的适应能力；但当窗口长度过长或 \(\gamma\) 取值过大时，虽能进一步降低 MAPE，但 AAS 却出现了下降的趋势。进一步从 Loss 曲线可见，各设置下损失均快速下降并趋近于稳定，表明训练过程整体可收敛；然而在小 \(\gamma\) 或短窗口长度条件下，MAPE 与 Loss、AAS 的收敛节奏并不同步，体现出特征学习在该条件下更易受噪声影响、稳定性不足。

综合预测准确性、模型鲁棒性与训练稳定性三者的权衡结果可以看出：当时间窗口长度为 72 且注意力参数为 2.0 时，MAPE、AAS 能够快速达到并稳定在较优水平，Loss 收敛平滑且无明显震荡。因此，该参数组合在预测精度、异常适应能力、训练稳定性三项指标之间表现出更优的均衡性。

\section{本章小结}\label{ux672cux7ae0ux5c0fux7ed3-1}

本章围绕传统时序预测方法在应对突发变化和数据噪声时的局限，提出了一种基于注意力机制的自适应时序预测方法。该方法构建时序特征与预测权重之间的映射机制：首先在时序域提取周期性、趋势性和波动性等特征，利用注意力机制形成时间级权重图并通过自适应融合为预测输入，以提升模型对关键信息的捕捉能力；随后在频域采用小波变换的多尺度特征提取，并结合自适应掩码与注意力参数调制，使特征表示更加符合时序数据的自然分布，从而在控制噪声影响的同时增强对异常变化和分布漂移的鲁棒性。在实验评估中，本章从预测准确性、模型鲁棒性、计算效率与泛化能力等维度构建完整评测体系，并通过对时间窗口长度 \(L\) 与注意力参数 \(\gamma\) 的消融实验，分析刻画预测精度与计算效率之间的权衡关系，验证了本章方法在保持预测精度的前提下，实现更优鲁棒性和更高计算效率的可行性与可控性。


% !TEX root = ../main.tex

\chapter{基于异常检测的故障诊断方法}\label{ch:system-management}

\section{引言}\label{ux5f15ux8a00-1}

异常检测与故障诊断（Anomaly Detection and Fault Diagnosis）目标是从运行数据中识别异常模式和定位故障根因，它是智能运维、系统可靠性和服务保障等应用的核心组件。在现代企业系统中，异常检测将参与故障预警、根因分析和系统恢复等关键环节。因此，系统的故障诊断能力对于保障业务连续性和降低运维成本具有重要意义。

与传统的固定规则检测相比，现代系统的故障模式具有三个显著特征：其一，故障类型多样且相互关联，单一故障可能引发连锁反应；其二，云原生架构下的微服务部署导致故障传播路径复杂，依赖关系难以追踪；其三，业务场景的动态变化使得预设的规则和阈值难以适应所有情况。这些因素共同导致传统的基于阈值的检测方法难以有效应对复杂系统的故障诊断；同时，现有的异常检测方法，如统计控制和机器学习分类等，也难以直接在海量、高维的监控数据中实现实时准确的故障定位。

基于上述挑战，本章面向复杂系统提出一种基于深度学习的异常检测方法。方法在特征设计上关注故障模式的有效表示，而不是简单地依赖人工设定的阈值。时序特征用于捕捉性能指标的演化模式，关联特征用于建模服务之间的依赖关系。两类特征均采用自适应提取和多尺度表示策略，使检测模型能够适应不同的业务场景和系统架构，从而提升检测的准确性和可解释性。

在检测目标方面，本章围绕异常诊断主题提出两类互补任务，点异常检测与群体异常检测。前者关注单个服务的异常行为，后者关注多个服务之间的关联异常。本章将异常检测过程形式化为特征提取与模式识别的联合优化，并给出严格的理论框架、符号定义、计算方法和评价指标体系；实验部分采用多个标准数据集和代表性场景，从检测精度、误报率、可解释性和实用性多个维度开展系统评估，整体框架如图~\ref{fig:ch4-system-monitoring-framework}~所示。

\begin{figure}[htbp]
\centering
\ThesisFigFillWidth{fig-ch4-01_image14.png}
\caption{基于深度学习的异常检测与故障诊断整体框架图}
\label{fig:ch4-system-monitoring-framework}
\end{figure}

\section{方案设计}\label{ux65b9ux6848ux8bbeux8ba1-1}

本节给出分布式系统监控的形式化定义与整体方案。分别给出监控数据采集和指标计算方法，最后定义两类互补的监控任务，并将其嵌入系统运维框架形成联合优化目标。

\subsection{问题与符号定义}\label{ux95eeux9898ux4e0eux7b26ux53f7ux5b9aux4e49}

（1）输入数据与系统状态。设输入数据由系统监控指标组成，任意时刻的指标向量记为 \(x_{t} \in \mathbb{R}^{d}\)，其中 \(d\) 是指标维度。时间窗口大小记为 \(T\)，目标时间点记为 \(t\)，相关时间序列记为 \(\mathcal{X}_{t} = \{x_{t-T}, \ldots, x_{t-1}, x_{t}\}\)。模型参数记为 \(\theta\)。

（2）异常检测机制。检测模型记为 \(f_{\theta}\)，输入为时间序列，输出为异常分数。本章统一记异常分数为 \(s_t = f_{\theta}(x_t, \mathcal{X}_{t-1})\)。关联矩阵记为 \(A \in \mathbb{R}^{d \times d}\)，其中 \(A_{ij}\) 表示指标 \(i\) 和指标 \(j\) 之间的相关性强度。

（3）自适应阈值。利用历史统计和动态更新等机制，系统可以根据数据分布变化自动调整检测阈值。对于点异常检测，系统需要根据历史分布确定正常范围的边界；对于群体异常检测，系统需要检测指标之间的相关性变化。同一指标记为 \(s_t\)；异常阈值记为 \(\tau\)，检测决策定义为：
\[
\hat{y}_t = \mathbb{I}(s_t > \tau)
\]
其中 \(\mathbb{I}(\cdot)\) 为指示函数，\(\hat{y}_t = 1\) 表示检测到异常。

\section{实验设置}\label{ux5b9eux9a8cux8bbeux7f6e}

\subsection{测试环境与系统}

为验证方法的有效性，我们在多个公开数据集和模拟环境中进行了实验。测试环境包含了服务器监控领域的多个数据集，如NASA服务器数据、KPI异常检测数据集和云服务监控数据等。基准环境作为测试环境的扩展，包含了更具挑战性的场景。在系统方面，我们选择了虚拟机集群、容器编排平台和无服务器架构三种代表性的部署环境。

\subsection{评价指标}

从检测精度、稳定性和可解释性三个维度进行评估。检测精度指标包括精确率、召回率和F1分数，衡量检测模型对异常的识别能力。稳定性指标包括不同噪声水平下的检测一致性和分布漂移下的鲁棒性，衡量检测方法的可靠性。可解释性指标包括特征重要性和决策路径分析，衡量检测结果的可信度和可理解性。

\section{实验结果与分析}

\subsection{有效性评估}

实验结果表明，我们的方法在各种数据集上都能取得良好的检测效果。在NASA服务器数据集上，检测模型能够准确识别CPU异常、内存泄漏等故障；在KPI异常检测数据集上，检测模型能够捕捉业务指标的异常波动；在云服务监控数据中，检测模型能够定位服务降级和响应延迟的原因。与基线方法相比，我们的方法在检测精度和误报率控制上都有显著提升。

\subsection{消融研究}

我们进行了详细的消融实验来验证各组件的作用。结果表明，时序特征提取对检测性能的贡献最大，能够有效捕捉性能指标的变化趋势；关联特征建模提供了依赖关系的表示能力，使检测更加准确；自适应阈值机制允许根据数据分布变化调整检测敏感度，从而适应不同场景的需求。

\subsection{案例分析}

通过具体的案例分析，我们展示了方法的实际应用效果。在服务器集群监控中，模型正确识别了导致服务降级的根本原因；在微服务架构中，模型准确标注了故障传播路径和受影响的服务；在容器编排平台中，异常分数变化显示了资源瓶颈的位置。这些案例验证了方法在实际应用中的有效性和实用性。
% !TEX root = ../main.tex

\chapter{总结与展望}\label{ch:conclusion}

\section{论文工作总结}\label{ux603bux7ed3}

本文围绕时序数据预测和异常检测问题展开研究，针对预测精度和故障诊断任务，分别提出了两类具有针对性的方法，并从预测准确性、模型鲁棒性、检测效果和实用性等方面进行了系统实验验证。研究结果表明，在本文设定的实验条件下，信息系统可通过智能预测方法和异常检测技术进一步提升预测准确性和故障定位能力。这说明智能系统在复杂时序场景下的预测和诊断能力仍值得持续关注。

（1）针对传统预测方法在面对突发变化和数据噪声时表现脆弱的问题，本文提出了一种基于注意力机制的智能预测方法。该方法综合利用时序特征提取、注意力机制和自适应融合策略，构建鲁棒的预测框架，以兼顾预测精度和模型鲁棒性。实验在多个公开数据集上开展，并结合预测误差、异常适应度、模型参数量和推理延迟等指标对方法性能进行综合评估。实验结果表明，所提方法在各种数据集上的预测误差均保持在较低水平，异常适应度均不低于0.8；在计算效率方面，模型参数量和推理延迟显著降低，说明该方法在时序预测方面具有较好的准确性和实用性。

（2）针对系统复杂性导致的故障定位困难问题，本文进一步提出了一种基于深度学习的异常检测方法。该方法面向复杂系统，设计了多层次的特征提取机制，构建了点异常检测和群体异常检测两类任务，使系统能够在保证检测精度的同时，提供可解释的故障诊断依据。实验在多个监控数据集上进行，并结合检测精度、误报率和可解释性等指标，从检测准确性、稳定性和实用性等方面进行了系统分析。实验结果表明，在保证检测精度的前提下，异常检测方法能够准确识别故障根因，提供符合运维直觉的诊断结果，用户满意度评分超过8.5（满分10分）。

需要指出的是，本文工作仍存在以下局限性：第一，面向时序预测的智能方法主要在单一任务场景上验证，在多任务学习场景上的表现仍需进一步评估；第二，面向异常检测的方法目前主要在静态数据集上验证，流式数据场景下的在线学习能力尚未系统考察；第三，两类方法均在理想实验环境下验证，在实际应用场景中的计算效率和实时性仍有待优化；第四，本文所采用的基线方法主要用于对比分析，针对特定领域的专用预测和检测方法仍需进一步研究。

\section{未来工作展望}\label{ux5c55ux671b}

尽管本文围绕时序预测和异常检测开展了初步研究，并取得了一定结果，但仍有若干问题值得在后续工作中进一步深入。

（1）面向实际应用场景的模型优化与部署。本文提出的方法主要在理想实验环境中验证。后续可进一步考虑不同数据分布、实时性要求和资源限制对模型性能的影响，研究更加高效和实用的模型优化和部署技术，从而提升在实际应用中的可用性。

（2）面向多模态和跨任务的预测方法。随着信息技术应用向多模态和跨任务扩展，单一预测模型的方法可能面临适用性不足的问题。未来有必要结合多模态数据的特性，进一步研究跨模态关联分析和统一预测框架，构建更具普适性的预测系统。

（3）异常检测评价体系的进一步完善。当前关于异常检测的评估仍主要围绕检测精度、误报率和用户满意度展开。对于多类型故障和复杂系统，如何进一步建立能够同时反映检测深度、覆盖范围、可解释性和实际诊断需求的综合评价指标体系，仍是值得持续研究的重要方向。

% !TEX root = ../main.tex
% 11 参考文献（排在正文章节之后；数据库见同目录 references.bib）
% 数字型著录顺序由正文首次引用顺序决定；2025 年文献在参考文献表中的前置由 main.tex 中 \begin{document} 后 \nocite{ref20,add02,add06,add07,ref29} 统一控制。
% \bibliographystyle{gbt7714-numerical}
\bibliographystyle{bst/gbt7714-numerical-local}
% 须先由 \bibsection 起章再起页，再 \addcontentsline，否则目录中「参考文献」会落在前一章末页
\renewcommand{\bibsection}{%
  \chapter*{\bibname}%
  \markboth{\bibname}{\bibname}%
  \phantomsection
  \addcontentsline{toc}{chapter}{\bibname}%
}
% 参考文献表版式：整词断行（行末不用连字符拆英文单词）、条目段落两端对齐（类 Word justified）
\begingroup
\justifying
\hyphenpenalty=10000\relax
\exhyphenpenalty=10000\relax
\sloppy
\emergencystretch=2.5em\relax
\bibliography{content/references}
\endgroup


% 科研成果、致谢：twoside 下从奇数页起排（必要时自动插入空白偶数页）
\cleardoublepage
% !TEX root = ../main.tex
%
% content/achievements上面文件夹下的材料，是我《攻读硕士学位期间的科研成果》的全部材料，分别是会议论文一篇，专利两项，请将这些材料整合到content/12-achievements.tex中，并进行适当的排版和格式(format.png是样例格式）调整，使内容更加清晰易读。编写tex代码的时候，请保留这些提示词作为注释内容。
% 要求：
% 1. 严谨准确，不得有任何错误；
% 2. 内容完整，不得有任何遗漏；
% 3. 本页标题、字体格式符合论文格式要求；
%
% 材料来源说明（便于核对）：
% - 会议论文：content/achievements/paper.pdf
% - 专利受理通知书：content/achievements/paper1.jpg、content/achievements/paper2.jpg
%
% 版式说明：本页按 content/achievements/format.png 参考样式排版——不出现本人姓名、单位；论文条目为「第一作者，英文题名，《出处》。（状态）」；专利为「申请发明专利，发明人角色说明，《名称》。（申请信息）」。
%
% 版式技术说明：minipage 内默认会 \@arrayparboxrestore 为两端对齐，\raggedright 易被覆盖；改用 ragged2e 的 \RaggedRight（按字体固定 \spaceskip）。
% 中英边界 xeCJK 的 CJKecglue 若含 plus 会在行末被拉伸，故在组内改为固定宽度；书名中 Computer Science 用 \mbox 防止词间被拆到两行。
%
% 奇数页起排：由 main.tex 在 \input 本文件前执行 \cleardoublepage

\chapter*{攻读硕士学位期间的科研成果}
\phantomsection
\addcontentsline{toc}{chapter}{攻读硕士学位期间的科研成果}
\thispagestyle{gzubody}
\markboth{攻读硕士学位期间的科研成果}{攻读硕士学位期间的科研成果}

\vspace{0.5\baselineskip}
\noindent
\begin{enumerate}[
  label={\textbf{[\arabic*]}},
  leftmargin=2.4em,
  labelsep=0.55em,
  itemsep=1.2\baselineskip,
  parsep=0pt,
  topsep=0pt,
  align=left
]

  \item
  \begingroup
  \xeCJKsetup{CJKecglue = \hskip 0.25em\relax}
  \begin{minipage}[t]{\linewidth}
    \RaggedRight
    \setlength{\parskip}{0pt}%
    \setlength{\parindent}{0pt}%
    {\songti 第一作者，}%
    {\rmfamily\normalfont A Study on Information System Performance Optimization and Reliability}%
    {\songti，《Lecture Notes in {\rmfamily\mbox{Computer Science}}》（International Conference on Information Systems），Springer，2025。（EI 收录，已发表）}%
  \end{minipage}
  \endgroup

  \item
  \begin{minipage}[t]{\linewidth}
    \RaggedRight
    \songti 申请发明专利，学生第四发明人，一种基于数据驱动的系统性能监控与优化方法。（申请号：202411059964.8；申请日：2024年8月5日；已受理）
  \end{minipage}

  \item
  \begin{minipage}[t]{\linewidth}
    \RaggedRight
    \songti 申请发明专利，学生第三发明人，一种面向分布式系统的资源调度与管理装置。（申请号：202410572817.4；申请日：2024年5月10日；已受理）
  \end{minipage}

\end{enumerate}

\cleardoublepage
% !TEX root = ../main.tex
% 致谢部分（正文为通用写法；定稿时请按实际补充导师姓名、课题组等具体致谢对象）
%
% 奇数页起排：由 main.tex 在 \input 本文件前执行 \cleardoublepage（必要时插入空白偶数页）
\chapter*{致谢}
\phantomsection
\addcontentsline{toc}{chapter}{致谢}
\thispagestyle{gzubody}
\markboth{致谢}{致谢}

攻读学位期间，本人在课程学习、科研工作与学位论文撰写过程中得到诸多师长、同窗及有关方面的帮助与支持，在此致以诚挚谢意。

感谢学校、学院和培养单位所提供的培养条件与学术氛围。感谢导师在论文选题、研究开展与撰稿修改等环节给予的指教。感谢课题组与实验室成员在学术交流、实验与工作条件等方面的协助。感谢评阅专家和答辩专家对论文提出的意见与建议。感谢家人与亲友在学途中的理解与鼓励。

囿于学识与能力所限，文中疏漏在所难免，敬请各位专家和读者指正。


\end{document}
